\cellcolor{alignmentblue}\textbf{B1 General Principles} & \cellcolor{alignmentblue}\textbf{10} & \cellcolor{alignmentblue}\textbf{30} & \cellcolor{alignmentblue}TBD \\
\rowcolor{alignmentblue}
B1.1 Biochemistry and molecular biology & & & \\
\rowcolor{alignmentblue}
B1.1.1 Structure, and function of proteins & & & \\
\rowcolor{alignmentblue}
B1.1.1.1 principles of protein structure and folding & & & \\
\rowcolor{alignmentblue}
B1.1.1.2 enzymes: kinetics, reaction mechanisms & & & \\
\rowcolor{alignmentblue}
B1.1.1.3 structural and regulatory proteins: ligand binding, self-assembly & & & \\
\rowcolor{alignmentblue}
B1.1.1.4 regulatory properties & & & \\
\rowcolor{alignmentblue}
B1.1.2 Gene expression: DNA structure, replication, and exchange & & & \\
\rowcolor{alignmentblue}
B1.1.2.1 DNA structure: single- and double-stranded DNA, stabilizing forces, supercoiling and topoisomers & & & \\
\rowcolor{alignmentblue}
B1.1.2.2 gene structure and organization: chromosomes, centromere, telomere & & & \\
\rowcolor{alignmentblue}
B1.1.2.3 DNA replication, degradation, repair, mutation, activation and inactivation & & & \\
\rowcolor{alignmentblue}
B1.1.2.4 recombination, insertion sequences, transposon & & & \\
\rowcolor{alignmentblue}
B1.1.2.5 mechanisms of genetic exchange, including transformation, transduction, conjugation, crossing over, recombination, linkage & & & \\
\rowcolor{alignmentblue}
B1.1.3 Gene expression: transcription, including defects & & & \\
\rowcolor{alignmentblue}
B1.1.3.1 transcription of DNA into RNA, RNA, RNA degradation & & & \\
\rowcolor{alignmentblue}
B1.1.3.2 post-transcriptional modification & & & \\
\rowcolor{alignmentblue}
B1.1.3.3 regulation: cis-regulatory elements, transcription factors, enhancers, promoters, silencers, repressants & & & \\
\rowcolor{alignmentblue}
B1.1.4 Gene expression: translation, including defects & & & \\
\rowcolor{alignmentblue}
B1.1.4.1 the genetic code & & & \\
\rowcolor{alignmentblue}
B1.1.4.2 structure and function of tRNA & & & \\
\rowcolor{alignmentblue}
B1.1.4.3 structure and function of ribosomes & & & \\
\rowcolor{alignmentblue}
B1.1.4.4 protein synthesis & & & \\
\rowcolor{alignmentblue}
B1.1.4.5 regulation of translation & & & \\
\rowcolor{alignmentblue}
B1.1.4.6 post-translational modifications & & & \\
\rowcolor{alignmentblue}
B1.1.5 Principles of molecular technology & & & \\
\rowcolor{alignmentblue}
B1.1.5.1 analysis of DNA: sequencing, restriction analysis, PCR, amplification, hybridization & & & \\
\rowcolor{alignmentblue}
B1.1.5.2 plasmids and bacteriophages & & & \\
\rowcolor{alignmentblue}
B1.1.5.3 gene cloning & & & \\
\rowcolor{alignmentblue}
B1.1.5.4 recombinant DNA technology and application & & & \\
\rowcolor{alignmentblue}
B1.1.6 Energy metabolism, including metabolic sequences, regulation & & & \\
\rowcolor{alignmentblue}
B1.1.6.1 cellular bioenergetic: thermodynamics, free energy, chemical equilibria and group transfer potential, energetics of ATP and other high-energy compounds & & & \\
\rowcolor{alignmentblue}
B1.1.6.2 tricarboxylic acid cycle, electron transport and oxidative phosphorylation & & & \\
\rowcolor{alignmentblue}
B1.1.6.3 generation of energy from carbohydrates, fatty acids, and essential amino acids; aerobic and anaerobic glycolysis and gluconeogenesis, β oxidation, ketogenesis and ketone bodies oxidation & & & \\
\rowcolor{alignmentblue}
B1.1.6.4 storage of energy: glycogenesis, fatty acid and triglyceride synthesis & & & \\
\rowcolor{alignmentblue}
B1.1.6.5 generation, expenditure, and storage of energy at the whole-body level & & & \\
\rowcolor{alignmentblue}
B1.1.6.6 protein turnover, protein-calorie balance & & & \\
\rowcolor{alignmentblue}
B1.1.7 Metabolic pathways of small molecules and associated diseases & & & \\
\rowcolor{alignmentblue}
B1.1.7.1 biosynthesis and degradation of carbohydrate including fructose, sorbitol, galactose, lactose, pentose phosphate pathway and uronic acid pathway (e.g., galactosemia, fructosuria, glucose 6-postphate dehydrogenase deficiency) & & & \\
\rowcolor{alignmentblue}
B1.1.7.2 biosynthesis and degradation of amino acids (e.g., phenylketonuria, homocystinuria, maple syrup urine diseases) & & & \\
\rowcolor{alignmentblue}
B1.1.7.3 biosynthesis and degradation of purine and pyrimidine nucleotides (e.g., gout, Lesch-Nyhan syndrome) & & & \\
\rowcolor{alignmentblue}
B1.1.7.4 biosynthesis and degradation of lipids (e.g., dyslipidemias, carnitine deficiency) & & & \\
\rowcolor{alignmentblue}
B1.1.7.5 biosynthesis and degradation of porphyrins (e.g., porphyria) & & & \\
\rowcolor{alignmentblue}
B1.1.7.6 biosynthesis and degradation of alcohols and other small molecules & & & \\
\rowcolor{alignmentblue}
B1.1.8 Biosynthesis and degradation of other macromolecules and associated abnormalities, complex carbohydrates (e.g., lysosomal storage disease), glycoproteins, and proteoglycans (e.g., type II glycogen storage disease) & & & \\
\rowcolor{alignmentblue}
B1.1.9 Nutrients, vitamins, minerals and trace elements & & & \\
\rowcolor{alignmentblue}
B1.1.9.1 functions of nutrients , including essential fatty acid, trans-fatty acids, cholesterol & & & \\
\rowcolor{alignmentblue}
B1.1.9.2 vitamins: structures of vitamins and derivatives of medical important, absorption, storage and excretion, functions, food sources and requirements & & & \\
\rowcolor{alignmentblue}
B1.1.9.3 mineral and trace element requirements & & & \\
\rowcolor{alignmentblue}
B1.1.9.4 assessment of nutritional status across the life span, including calories, protein, essential nutrients, hypoalimentation & & & \\
\rowcolor{alignmentblue}
B1.2 Biology of cells & & & \\
\rowcolor{alignmentblue}
B1.2.1 Structure and function of cell components (e.g., endoplasmic reticulum, Golgi complex, mitochondria, lysosome, peroxisome, endosome, centriole, microtubule, ribosome, polysome, plasma membrane, cytosol, cilia, nucleus, cytoskeleton) & & & \\
\rowcolor{alignmentblue}
B1.2.2 Signal transduction (including basic principles, receptors and channels, second messengers, signal transduction pathways) & & & \\
\rowcolor{alignmentblue}
B1.2.3 Cell-cell and cell-matrix adhesion & & & \\
\rowcolor{alignmentblue}
B1.2.4 Cell motility & & & \\
\rowcolor{alignmentblue}
B1.2.5 Intracellular sorting (e.g., trafficking, endocytosis) & & & \\
\rowcolor{alignmentblue}
B1.2.6 Cellular homeostasis (e.g., turnover, pH maintenance, proteasome, ions, soluble proteins) & & & \\
\rowcolor{alignmentblue}
B1.2.7 Cell cycle (e.g., mitosis, meiosis, structure of spindle apparatus, cell cycle regulation) & & & \\
\rowcolor{alignmentblue}
B1.2.8 Structure and function of basic tissue components (including epithelial cells, connective tissue cells, muscle cells, nerve cells, and extracellular matrix) & & & \\
\rowcolor{alignmentblue}
B1.2.9 Intracellular accumulations (e.g., pigments, fats, proteins, carbohydrates, minerals, Inclusions, vacuoles) & & & \\
\rowcolor{alignmentblue}
B1.2.10 Adaptive cell response to injury (e.g., hypertrophy and metaplasia) & & & \\
\rowcolor{alignmentblue}
B1.2.11 Mechanisms of cell injury and necrosis, including free radical mediated cell injury & & & \\
\rowcolor{alignmentblue}
B1.2.12 Apoptosis & & & \\
\rowcolor{alignmentblue}
B1.3 Human development and genetics & & & \\
\rowcolor{alignmentblue}
B1.3.1 Embryogenesis: basic programmed gene expression, and developmental regulation of gene expression for medical student & & & \\
\rowcolor{alignmentblue}
B1.3.2 Principles of pedigree analysis, including inheritance patterns, occurrence and recurrence risk determination & & & \\
\rowcolor{alignmentblue}
B1.3.3 Population genetics: Hardy-Weinberg law, founder effects, mutation-selection equilibrium & & & \\
\rowcolor{alignmentblue}
B1.4 Normal immune responses & & & \\
\rowcolor{alignmentblue}
B1.4.1 Innate immunity & & & \\
\rowcolor{alignmentblue}
B1.4.2 Production and function of granulocytes, natural killer cells, macrophages & & & \\
\rowcolor{alignmentblue}
B1.4.3 Production and function of T lymphocytes, T-lymphocyte receptors & & & \\
\rowcolor{alignmentblue}
B1.4.4 Production and function of B lymphocytes and plasma cells, immunoglobulin and antibodies, structure and biologic properties & & & \\
\rowcolor{alignmentblue}
B1.4.5 Antigenicity and immunogenicity, antigen presentation, cell activation and regulation, tolerance and clonal deletion & & & \\
\rowcolor{alignmentblue}
B1.4.6 Immunologic mediators: chemistry, function, molecular biology, complement pathways, cytokines, chemokines & & & \\
\rowcolor{alignmentblue}
B1.4.7 Immunogenetics: MHC structure and function, erythrocyte antigens, transplantation & & & \\
\rowcolor{alignmentblue}
B1.4.8 Immunizations, protective immunity & & & \\
\rowcolor{alignmentblue}
B1.5 Pathogenesis, pathophysiology, basic pathological process and laboratory investigation & & & \\
\rowcolor{alignmentblue}
B1.5.1 Acute and chronic inflammation & & & \\
\rowcolor{alignmentblue}
B1.5.1.1 acute inflammation and mediator systems & & & \\
\rowcolor{alignmentblue}
B1.5.1.2 bactericidal mechanisms and tissue injury & & & \\
\rowcolor{alignmentblue}
B1.5.1.3 clinical manifestations (e.g., pain, leukocytosis, fever, leukemoid reaction, and chills) & & & \\
\rowcolor{alignmentblue}
B1.5.1.4 chronic inflammation & & & \\
\rowcolor{alignmentblue}
B1.5.2 Tissue renewal, regeneration and repair & & & \\
\rowcolor{alignmentblue}
B1.5.2.1 control of normal cell proliferation and tissue growth (e.g., stem cell, tissue proliferative activity) & & & \\
\rowcolor{alignmentblue}
B1.5.2.2 mechanism of tissue and organ regeneration & & & \\
\rowcolor{alignmentblue}
B1.5.2.3 healing by repair, scar formation and fibrosis & & & \\
\rowcolor{alignmentblue}
B1.5.3 Hemodynamic disorders, thromboembolic disease and shock & & & \\
\rowcolor{alignmentblue}
B1.5.3.1 edema & & & \\
\rowcolor{alignmentblue}
B1.5.3.2 hyperemia and congestion & & & \\
\rowcolor{alignmentblue}
B1.5.3.3 hemorrhage & & & \\
\rowcolor{alignmentblue}
B1.5.3.4 hemostasis and thrombosis & & & \\
\rowcolor{alignmentblue}
B1.5.3.5 embolism (e.g., pulmonary, air, fat, bone marrow, amniotic embolism, systemic thromboembolism) & & & \\
\rowcolor{alignmentblue}
B1.5.3.6 infarction & & & \\
\rowcolor{alignmentblue}
B1.5.3.7 shock & & & \\
\rowcolor{alignmentblue}
B1.5.4 Genetic disorders & & & \\
\rowcolor{alignmentblue}
B1.5.4.1 gene and human diseases (e.g., mutations) & & & \\
\rowcolor{alignmentblue}
B1.5.4.2 congenital abnormalities: principles, patterns of anomalies, dysmorphogenesis & & & \\
\rowcolor{alignmentblue}
B1.5.4.3 genetic mechanisms: chromosomal abnormalities, Mendelian inheritance, complex multigenic disorders, single gene disorders with nonclassic inheritance & & & \\
\rowcolor{alignmentblue}
B1.5.4.3 clinical genetics, including genetic testing, prenatal diagnosis, newborn screening, genetic counseling/ethics, gene therapy & & & \\
\rowcolor{alignmentblue}
B1.5.5 Diseases of the immune system & & & \\
\rowcolor{alignmentblue}
B1.5.5.1 immunodeficiency & & & \\
\rowcolor{alignmentblue}
B1.5.5.1.1 primary Immunodeficiency & & & \\
\rowcolor{alignmentblue}
B1.5.5.1.1.1 combined immunodeficiency disease (e.g., SCID) & & & \\
\rowcolor{alignmentblue}
B1.5.5.1.1.2 T cell deficiency (e.g., DiGeorge syndrome) & & & \\
\rowcolor{alignmentblue}
B1.5.5.1.1.3 B cell deficiency (e.g., congenital agammaglobulinemia) & & & \\
\rowcolor{alignmentblue}
B1.5.5.1.1.4 phagocyte deficiency (e.g., chronic granulomatous disease) & & & \\
\rowcolor{alignmentblue}
B1.5.5.1.1.5 complement deficiency & & & \\
\rowcolor{alignmentblue}
B1.5.5.1.2 secondary Immunodeficiency (e.g., infection, radiation, splenectomy, malnutrition etc.) & & & \\
\rowcolor{alignmentblue}
B1.5.5.2 immunologically mediated disorders & & & \\
\rowcolor{alignmentblue}
B1.5.5.2.1 hypersensitivity types I–IV & & & \\
\rowcolor{alignmentblue}
B1.5.5.2.2 autoimmune disorders & & & \\
\rowcolor{alignmentblue}
B1.5.5.2.3 transplantation and graft rejection including graft VS host reaction & & & \\
\rowcolor{alignmentblue}
B1.5.5.3 tumor immunology & & & \\
\rowcolor{alignmentblue}
B1.5.5.4 laboratory investigation (e.g., ELISA, complement fixation RIA, agglutination) & & & \\
\rowcolor{alignmentblue}
B1.5.5.5 principle of immunotherapy & & & \\
\rowcolor{alignmentblue}
B1.5.6 Neoplasia & & & \\
\rowcolor{alignmentblue}
B1.5.6.1 classification, histologic diagnosis & & & \\
\rowcolor{alignmentblue}
B1.5.6.2 grading and staging of neoplasms & & & \\
\rowcolor{alignmentblue}
B1.5.6.3 cell biology, biochemistry, and molecular biology of neoplastic cells: transformation, oncogenes, oncogenesis, altered cell differentiation, and proliferation & & & \\
\rowcolor{alignmentblue}
B1.5.6.4 molecular basis of multistep carcinogenesis & & & \\
\rowcolor{alignmentblue}
B1.5.6.4 carcinogenic agents (e.g., chemical, radiation, and microbial agents) & & & \\
\rowcolor{alignmentblue}
B1.5.6.5 invasion and metastasis & & & \\
\rowcolor{alignmentblue}
B1.5.6.6 host defense against tumors & & & \\
\rowcolor{alignmentblue}
B1.5.6.7 paraneoplastic manifestations of cancer & & & \\
\rowcolor{alignmentblue}
B1.5.6.8 cancer epidemiology and prevention & & & \\
\rowcolor{alignmentblue}
B1.5.7 Infectious diseases & & & \\
\rowcolor{alignmentblue}
B1.5.7.1 microbial classification & & & \\
\rowcolor{alignmentblue}
B1.5.7.2 bacteria and bacterial diseases & & & \\
\rowcolor{alignmentblue}
B1.5.7.2.1 structure and composition & & & \\
\rowcolor{alignmentblue}
B1.5.7.2.2 metabolism, physiology, and regulation & & & \\
\rowcolor{alignmentblue}
B1.5.7.2.3 genetics & & & \\
\rowcolor{alignmentblue}
B1.5.7.2.4 nature and mechanisms of action of virulence factors & & & \\
\rowcolor{alignmentblue}
B1.5.7.2.5 pathophysiology of infection & & & \\
\rowcolor{alignmentblue}
B1.5.7.2.6 epidemiology and ecology & & & \\
\rowcolor{alignmentblue}
B1.5.7.2.7 principles of cultivation, assay, and laboratory diagnosis & & & \\
\rowcolor{alignmentblue}
B1.5.7.2.8 bacterial infection not classified in the system; melioidosis, leptospirosis, diphtheria, typhus fever (scrub typhus, murine typhus), pertussis, anthrax, meningococcemia & & & \\
\rowcolor{alignmentblue}
B1.5.7.3 viruses and viral diseases & & & \\
\rowcolor{alignmentblue}
B1.5.7.3.1 physical and chemical properties & & & \\
\rowcolor{alignmentblue}
B1.5.7.3.2 replication & & & \\
\rowcolor{alignmentblue}
B1.5.7.3.3 genetics & & & \\
\rowcolor{alignmentblue}
B1.5.7.3.4 principles of cultivation, assay, and laboratory diagnosis & & & \\
\rowcolor{alignmentblue}
B1.5.7.3.5 molecular basis of pathogenesis & & & \\
\rowcolor{alignmentblue}
B1.5.7.3.6 pathophysiology of infection & & & \\
\rowcolor{alignmentblue}
B1.5.7.3.7 latent and persistent infections & & & \\
\rowcolor{alignmentblue}
B1.5.7.3.8 epidemiology & & & \\
\rowcolor{alignmentblue}
B1.5.7.3.9 oncogenic viruses & & & \\
\rowcolor{alignmentblue}
B1.5.7.3.10 viral diseases not classified in the system; Dengue hemorrhagic fever \& Chikungunya fever, varicella, zoster, herpes simplex infection viral exanthemata (measles, rubella, erythema infectiosum, exanthema subitum, hand foot mouth diseases), mumps, infectious mononucleosis, cytomegaloviral infection & & & \\
\rowcolor{alignmentblue}
B1.5.7.4 fungi and fungal infections & & & \\
\rowcolor{alignmentblue}
B1.5.7.4.1 structure, physiology, cultivation, and laboratory diagnosis & & & \\
\rowcolor{alignmentblue}
B1.5.7.4.2 pathogenesis and epidemiology & & & \\
\rowcolor{alignmentblue}
B1.5.7.4.3 fungi and fungal infections not classified in the system; systemic mycoses (candidiasis, cryptococcosis, aspergillosis, penicillosis, mucormycosis) & & & \\
\rowcolor{alignmentblue}
B1.5.7.5 parasites and parasitic diseases & & & \\
\rowcolor{alignmentblue}
B1.5.7.5.1 structure, physiology, life cycles, and laboratory diagnosis & & & \\
\rowcolor{alignmentblue}
B1.5.7.5.2 pathogenesis and epidemiology & & & \\
\rowcolor{alignmentblue}
B1.5.7.5.3 parasitic diseases not classified in the system; blood and tissue parasitic diseases (trichinosis, cysticercosis, gnathostomiasis, pneumocystosis, angiostrongyliasis, filariasis, toxoplasmosis, opisthorchiasis, paragonimiasis, schistosomiasis, malaria, acanthamoeba infection, naegleria infection, leishmaniasis) & & & \\
\rowcolor{alignmentblue}
B1.5.7.6 principles of sterilization and pure culture technique & & & \\
\rowcolor{alignmentblue}
B1.5.8 Environmental and nutritional diseases & & & \\
\rowcolor{alignmentblue}
B1.5.8.1 adaptation to environmental extremes, including occupational exposures & & & \\
\rowcolor{alignmentblue}
B1.5.8.1.1 physical and associated disorders (e.g., mechanical trauma, radiation, decreased atmospheric pressure, thermal injuries, electrical injury, high- altitude sickness, increased water pressure) & & & \\
\rowcolor{alignmentblue}
B1.5.8.1.2 chemical poisoning and therapy (e.g., agricultural hazards, biological venom and toxins, heavy metal, gases, vapors, smoke inhalation, agricultural hazards, volatile organic solvents, principles of poisoning and therapy) & & & \\
\rowcolor{alignmentblue}
B1.5.8.2 nutritional disorders & & & \\
\rowcolor{alignmentblue}
B1.5.8.2.1 protein-calorie malnutrition & & & \\
\rowcolor{alignmentblue}
B1.5.8.2.2 vitamin deficiencies and toxicities & & & \\
\rowcolor{alignmentblue}
B1.5.8.2.3 mineral and trace element deficiencies and toxicities & & & \\
\rowcolor{alignmentblue}
B1.5.8.2.4 obesity & & & \\
\rowcolor{alignmentblue}
B1.5.8.2.5 alternative diets and nutritional supplements & & & \\
\rowcolor{alignmentblue}
B1.5.9 Principles of specimen collections and laboratory interpretations & & & \\
\rowcolor{alignmentblue}
B1.5.9.1 complete blood count; hemoglobin, hematocrit, white blood count, platelet count, red cell indices, blood smear: differential white blood cell count, red cell morphology, platelet estimate and morphology & & & \\
\rowcolor{alignmentblue}
B1.5.9.2 reticulocyte count & & & \\
\rowcolor{alignmentblue}
B1.5.9.3 hemoglobin typing & & & \\
\rowcolor{alignmentblue}
B1.5.9.4 inclusion bodies, Heinz bodies & & & \\
\rowcolor{alignmentblue}
B1.5.9.5 venous clotting time, clot retraction, clot lysis & & & \\
\rowcolor{alignmentblue}
B1.5.9.6 coagulation study & & & \\
\rowcolor{alignmentblue}
B1.5.9.7 bleeding time & & & \\
\rowcolor{alignmentblue}
B1.5.9.8 erythrocyte sedimentation rate & & & \\
\rowcolor{alignmentblue}
B1.5.9.9 malarial parasite; thick and thin film & & & \\
\rowcolor{alignmentblue}
B1.5.9.10 ABO blood group, Rh and cross matching & & & \\
\rowcolor{alignmentblue}
B1.5.9.11 urinalysis & & & \\
\rowcolor{alignmentblue}
B1.5.9.12 stool examination, stool occult blood & & & \\
\rowcolor{alignmentblue}
B1.5.9.13 Gram staining, acid fast staining & & & \\
\rowcolor{alignmentblue}
B1.5.9.14 KOH smear, Tzanck’s smear, wet preparation & & & \\
\rowcolor{alignmentblue}
B1.5.9.15 body fluid analysis (e.g., CSF, pleural) & & & \\
\rowcolor{alignmentblue}
B1.5.9.16 cytology & & & \\
\rowcolor{alignmentblue}
B1.5.9.17 culture from clinical specimens & & & \\
\rowcolor{alignmentblue}
B1.5.9.18 antimicrobial susceptibility test & & & \\
\rowcolor{alignmentblue}
B1.5.9.19 endocrinologic studies: plasma glucose, HbA1C, oral glucose tolerance test (OGTT), postprandial GTT, thyroid function test, serum cortisol & & & \\
\rowcolor{alignmentblue}
B1.5.9.20 liver profile & & & \\
\rowcolor{alignmentblue}
B1.5.9.21 kidney profile & & & \\
\rowcolor{alignmentblue}
B1.5.9.22 lipid profile & & & \\
\rowcolor{alignmentblue}
B1.5.9.23 cardiac markers & & & \\
\rowcolor{alignmentblue}
B1.5.9.24 tumor markers & & & \\
\rowcolor{alignmentblue}
B1.5.9.25 HIV testing & & & \\
\rowcolor{alignmentblue}
B1.5.9.26 serologic studies & & & \\
\rowcolor{alignmentblue}
B1.5.9.27 arterial blood gas analysis & & & \\
\rowcolor{alignmentblue}
B1.5.9.28 electrocardiography & & & \\
\rowcolor{alignmentblue}
B1.5.9.29 spirometry, peak expiratory flow rate measurement & & & \\
\rowcolor{alignmentblue}
B1.6 Gender, ethnic, and behavioral considerations affecting disease treatment and prevention, including psychosocial, cultural, occupational, and environmental & & & \\
\rowcolor{alignmentblue}
B1.6.1 Progression through the life cycle, including birth through senescence & & & \\
\rowcolor{alignmentblue}
B1.6.1.1 normal growth and development through each period & & & \\
\rowcolor{alignmentblue}
B1.6.1.2 psychosocial development of infancy, toddlers, preschool, school age, adolescence, young adults, middle age, menopause, aging male and old age & & & \\
\rowcolor{alignmentblue}
B1.6.1.3 cognitive, language, motor skills, social and interpersonal development & & & \\
\rowcolor{alignmentblue}
B1.6.1.4 sexual development (e.g., puberty, menopause) & & & \\
\rowcolor{alignmentblue}
B1.6.1.5 influence of developmental stage on physician/patient interview & & & \\
\rowcolor{alignmentblue}
B1.6.1.6 changes related to pregnancy & & & \\
\rowcolor{alignmentblue}
B1.6.1.7 personality development & & & \\
\rowcolor{alignmentblue}
B1.6.1.8 physical performance across lifespan & & & \\
\rowcolor{alignmentblue}
B1.6.1.9 sex and gender & & & \\
\rowcolor{alignmentblue}
B1.6.1.9.1 gender equity, gender equality, gender discrimination and social construction of gender & & & \\
\rowcolor{alignmentblue}
B1.6.1.9.2 man and masculinity & & & \\
\rowcolor{alignmentblue}
B1.6.1.9.3 gender and health (link between gender and health) & & & \\
\rowcolor{alignmentblue}
B1.6.1.9.4 social determinants of health: gender & & & \\
\rowcolor{alignmentblue}
B1.6.1.9.5 gender analysis, gender bias in medical education and in health care setting & & & \\
\rowcolor{alignmentblue}
B1.6.1.9.6 human rights & & & \\
\rowcolor{alignmentblue}
B1.6.1.9.7 gender sensitivity & & & \\
\rowcolor{alignmentblue}
B1.6.2 Psychologic and social factors influencing patient behavior & & & \\
\rowcolor{alignmentblue}
B1.6.2.1 personality traits or coping style, including coping mechanisms & & & \\
\rowcolor{alignmentblue}
B1.6.2.2 psychodynamic and behavioral factors, related past experience & & & \\
\rowcolor{alignmentblue}
B1.6.2.3 family and cultural factors, including socioeconomic status, ethnicity, and gender & & & \\
\rowcolor{alignmentblue}
B1.6.2.4 adaptive and maladaptive behavioral responses to stress and illness (e.g., drug- seeking behavior, sleep deprivation) & & & \\
\rowcolor{alignmentblue}
B1.6.2.5 interactions between the patient and the physician or the health care system (e.g., patient transfer) & & & \\
\rowcolor{alignmentblue}
B1.6.2.6 patient adherence, including general and adolescent & & & \\
\rowcolor{alignmentblue}
B1.6.3 Patient interviewing, consultation, and interactions with the family & & & \\
\rowcolor{alignmentblue}
B1.6.3.1 establishing and maintaining rapport & & & \\
\rowcolor{alignmentblue}
B1.6.3.2 data gathering & & & \\
\rowcolor{alignmentblue}
B1.6.3.3 approaches to patient education including patient’s understanding about the health problems, purpose of therapy, proper use of drugs, adherence to therapy and signs of major adverse drug effects & & & \\
\rowcolor{alignmentblue}
B1.6.3.4 enticing patients to make lifestyle changes & & & \\
\rowcolor{alignmentblue}
B1.6.3.5 communicating bad news & & & \\
\rowcolor{alignmentblue}
B1.6.3.6 “difficult” interviews (e.g., anxious or angry patients) & & & \\
\rowcolor{alignmentblue}
B1.6.3.7 multicultural ethnic characteristics & & & \\
\rowcolor{alignmentblue}
B1.6.4 Medical ethics, jurisprudence, and professional behavior & & & \\
\rowcolor{alignmentblue}
B1.6.4.1 principles of medical ethics & & & \\
\rowcolor{alignmentblue}
B1.6.4.2 doctor-patient relationship & & & \\
\rowcolor{alignmentblue}
B1.6.4.3 codes of conduct and etiquette & & & \\
\rowcolor{alignmentblue}
B1.6.4.4 human rights & & & \\
\rowcolor{alignmentblue}
B1.6.4.5 patient autonomy & & & \\
\rowcolor{alignmentblue}
B1.6.4.6 equity and social justice & & & \\
\rowcolor{alignmentblue}
B1.6.4.7 relationship with the pharmaceutical industry & & & \\
\rowcolor{alignmentblue}
B1.6.4.8 professionalism & & & \\
\rowcolor{alignmentblue}
B1.6.4.9 ethical dilemmas facing medical students & & & \\
\rowcolor{alignmentblue}
B1.6.4.10 principles of research ethics and research publication & & & \\
\rowcolor{alignmentblue}
B1.6.4.11 ethical aspects of genetics & & & \\
\rowcolor{alignmentblue}
B1.6.4.12 ethics of public health and health promotion & & & \\
\rowcolor{alignmentblue}
B1.6.5 Health schemes, national drug policies & & & \\
\rowcolor{alignmentblue}
B1.6.6 Organization and cost of health care delivery & & & \\
\rowcolor{alignmentblue}
B1.6.7 Concept and utilization of national list of essential medicines (NLEM) & & & \\
\rowcolor{alignmentblue}
B1.7 Multisystem processes & & & \\
\rowcolor{alignmentblue}
B1.7.1 Temperature regulation & & & \\
\rowcolor{alignmentblue}
B1.7.2 Fluid, electrolyte, and acid-base balance and disorders (e.g., water excess, dehydration, acidosis, alkalosis, hyponatremia, hypernatremia, hypokalemia, hyperkalemia, hypocalcemia, hypercalcemia) & & & \\
\rowcolor{alignmentblue}
B1.7.3 Exercise physiology & & & \\
\rowcolor{alignmentblue}
B1.7.4 Sleep health and sleep physiology & & & \\
\rowcolor{alignmentblue}
B1.8 General pharmacology & & & \\
\rowcolor{alignmentblue}
B1.8.1 Pharmacodynamic and pharmacokinetic processes & & & \\
\rowcolor{alignmentblue}
B1.8.1.1 pharmacokinetics: absorption and bioavailability, distribution, metabolism, excretion, dose, method of administration and dosage intervals & & & \\
\rowcolor{alignmentblue}
B1.8.1.2 pharmacodynamics, mechanisms of drug action, structure-activity relationships, receptors,signal transduction & & & \\
\rowcolor{alignmentblue}
B1.8.1.3 concentration- and dose-effect relationships (e.g., efficacy, potency), types and actions of agonists and antagonists & & & \\
\rowcolor{alignmentblue}
B1.8.1.4 individual factors or special populations altering pharmacokinetics and pharmacodynamics (e.g., age, pregnancy, gender, disease, tolerance, compliance, body weight, metabolic proficiency, pharmacogenetics) & & & \\
\rowcolor{alignmentblue}
B1.8.1.5 adverse drug effects including contraindication, warning, precautions, side effects, drug interactions, overdose and toxicity & & & \\
\rowcolor{alignmentblue}
B1.8.1.6 drug development and regulatory issues (e.g., approval scheduling, bioequivalence, biosimilar) & & & \\
\rowcolor{alignmentblue}
B1.8.1.7 pharmacoeconomics and pharmacoepidemiology & & & \\
\rowcolor{alignmentblue}
B1.8.1.8 phamarcovigilance (e.g., ADR reporting system) & & & \\
\rowcolor{alignmentblue}
B1.8.1.9 concept of rational drug prescription & & & \\
\rowcolor{alignmentblue}
B1.8.2 General properties of autacoids, including peptides and analogs, biogenic amines, prostanoids and their inhibitors, and smooth muscle/endothelial autacoids & & & \\
\rowcolor{alignmentblue}
B1.8.3 General principles of autonomic pharmacology & & & \\
\rowcolor{alignmentblue}
B1.8.4 General properties of antimicrobials, antiparasitic agents including mechanisms of action ; and mechanisms, prevention and management of drug resistance & & & \\
\rowcolor{alignmentblue}
B1.8.5 General properties of antineoplastic agents and immunosuppressants, including drug effects on rapidly dividing mammalian cells & & & \\
\rowcolor{alignmentblue}
B1.9 Quantitative methods & & & \\
\rowcolor{alignmentblue}
B1.9.1 Biostatistics & & & \\
\rowcolor{alignmentblue}
B1.9.1.1 statistical concepts in medical practice & & & \\
\rowcolor{alignmentblue}
B1.9.1.2 quantitative and qualitative data & & & \\
\rowcolor{alignmentblue}
B1.9.1.3 types of variables & & & \\
\rowcolor{alignmentblue}
B1.9.1.4 frequency distributions & & & \\
\rowcolor{alignmentblue}
B1.9.1.5 measures of central tendency & & & \\
\rowcolor{alignmentblue}
B1.9.1.6 measures of dispersion & & & \\
\rowcolor{alignmentblue}
B1.9.1.7 disease frequency, rates, risk & & & \\
\rowcolor{alignmentblue}
B1.9.1.8 statistical inference & & & \\
\rowcolor{alignmentblue}
B1.9.1.9 hypothesis generation, hypothesis testing, & & & \\
\rowcolor{alignmentblue}
B1.9.1.10 the alpha level and p value, & & & \\
\rowcolor{alignmentblue}
B1.9.1.11 type I error and type II error & & & \\
\rowcolor{alignmentblue}
B1.9.2 Research Methodology & & & \\
\rowcolor{alignmentblue}
B1.9.2.1 appraisal and application of medical literature & & & \\
\rowcolor{alignmentblue}
B1.9.2.2 foreground and background question, PICO & & & \\
\rowcolor{alignmentblue}
B1.9.2.3 sources and hierarchy of evidence & & & \\
\rowcolor{alignmentblue}
B1.9.2.4 research design & & & \\
\rowcolor{alignmentblue}
B1.9.2.5 study population, sampling, sample size & & & \\
\rowcolor{alignmentblue}
B1.9.2.6 sample allocation, concealment & & & \\
\rowcolor{alignmentblue}
B1.9.2.7 outcome assessment & & & \\
\rowcolor{alignmentblue}
B1.9.2.8 test of statistical significance, clinical important & & & \\
\rowcolor{alignmentblue}
B1.9.2.9 confidence intervals & & & \\
\rowcolor{alignmentblue}
B1.9.2.10 internal validity, external validity, reliability & & & \\
\rowcolor{alignmentblue}
B1.9.2.11 impact (size of effect) & & & \\
\rowcolor{alignmentblue}
B1.9.2.12 applicability & & & \\
\rowcolor{alignmentblue}
B1.9.2.13 patient values, patient circumstance & & & \\
\rowcolor{alignmentblue}
B1.9.2.14 chance, bias, confounder & & & \\
\rowcolor{alignmentblue}
B1.9.2.15 relative risk, odds ratio & & & \\
\rowcolor{alignmentblue}
B1.9.2.16 sensitivity, specificity, predictive value, likelihood ratio, ROC curve & & & \\
\rowcolor{alignmentblue}
B1.9.2.17 compliance & & & \\
\rowcolor{alignmentblue}
B1.9.2.18 co-intervention, contamination & & & \\
\rowcolor{alignmentblue}
B1.9.2.19 intention to treat, per-protocol analysis & & & \\
\rowcolor{alignmentblue}
B1.9.2.20 target (ultimate), surrogate (intermediate, substitute), outcomes (endpoints) & & & \\
\rowcolor{alignmentblue}
B1.9.2.21 relative risk reduction, absolute risk reduction, NNT & & & \\
\rowcolor{alignmentblue}
B1.9.2.22 cause and effect association หมวดที่ 2. การจ าแนกตามระบบอวัยวะ & & & \\
\rowcolor{alignmentblue}
ในหมวดนี้ได้ระบุเนื้อหาโดยจ าแนกตามระบบอวัยวะออกเป็น 10 ระบบ (B2-B11) ในแต่ละระบบประกอบด้วย 2.1 ความรู้เรื่องภาวะปกติ 2.2 ความรู้เรื่องสาเหตุ พยาธิก าเนิด พยาธิสรีรวิทยา และพยาธิสภาพ ตลอดจนหลักการวินิจฉัยสาเหตุ ของการ เกิด โรค ซึ่งแบ่งเป็น กลุ่มโรคและกลุ่มอาการที่ต้องรู้ และควรรู้โดยอิงจากเกณฑ์ในภาคผนวก ข ความรู้ ความสามารถทางวิชาชีพ และทักษะทางคลินิกดังนี้ - กลุ่มโรคและกลุ่มอาการที่ต้องรู้ อิงจากกลุ่มที่ 1 และ 2 รวมทั้งเนื้อหาอื่นๆ - กลุ่มโรคและกลุ่มอาการที่ควรรู้ อิงจากกลุ่มที่ 3 รวมทั้งเนื้อหา อื่นๆ 2.3 หลักการรักษาและหลักการใช้ยา โดยเฉพาะยาในบัญชียาหลักแห่งชาติ 2.4 การตระหนักถึงความส าคัญของพฤติกรรมสุขภาพ เพศ วัย เชื้อชาติ อาชีพ เศรษฐกิจ สังคม วัฒนธรรมและ สิ่งแวดล้อม ตลอดจนปัจจัยอื่นๆ ที่มีผลกระทบต่อการบ าบัดรักษา ฟื้นฟูสภาพ สร้างเสริมสุขภาพและ ป้องกันโรค & & & \\
\rowcolor{alignmentblue}
\textbf{B2 Hematopoietic and Lymphoreticular Systems} & \textbf{10} & \textbf{30} & TBD \\
\rowcolor{alignmentblue}
B2.1 มีความรู้เรื่องภาวะปกติของระบบโลหิตวิทยา ดังต่อไปนี้ & & & \\
\rowcolor{alignmentblue}
B2.1.1 Embryonic development, fetal maturation, and perinatal changes & & & \\
\rowcolor{alignmentblue}
B2.1.2 Organ structure and function (bone marrow, lymph node, thymus, spleen) & & & \\
\rowcolor{alignmentblue}
B2.1.3 Cell/tissue structure and function & & & \\
\rowcolor{alignmentblue}
B2.1.3.1 production and function of erythrocytes, hemoglobin, O2 and CO2 transport, transport proteins & & & \\
\rowcolor{alignmentblue}
B2.1.3.2 production and function of leukocytes and the lymphoreticular system & & & \\
\rowcolor{alignmentblue}
B2.1.3.3 production and function of platelets & & & \\
\rowcolor{alignmentblue}
B2.1.3.4 production and function of coagulation, fibrinolytic factors and natural inhibitors & & & \\
\rowcolor{alignmentblue}
B2.1.4. Repair, regeneration, and changes associated with stage of life & & & \\
\rowcolor{alignmentblue}
B2.2 มีความรู้เรื่องสาเหตุ พยาธิก าเนิด พยาธิสรีรวิทยา และพยาธิสภาพ ตลอดจนหลักการวินิจฉัยสาเหตุ ของ การเกิดโรค กลุ่มอาการและภาวะผิดปกติทางระบบโลหิตวิทยา ดังต่อไปนี้ & & & \\
\rowcolor{alignmentblue}
B2.2.1 Infectious, inflammatory, and immunologic disorders & & & \\
\rowcolor{alignmentyellow}
B2.2.1.1 กลุ่มโรคและกลุ่มอาการที่ต้องรู้ & & & \\
\rowcolor{alignmentred}
B2.2.1.2 กลุ่มโรคและกลุ่มอาการที่ควรรู้ & & & \\
\rowcolor{alignmentyellow}
(1) lymphoid hyperplasia & & & \\
\rowcolor{alignmentred}
(1) autoimmune hemolytic anemia & & & \\
\rowcolor{alignmentyellow}
(2) incompatible blood transfusion & & & \\
\rowcolor{alignmentred}
(2) immune thrombocytopenia & & & \\
\rowcolor{alignmentred}
(3) immune deficiency states & & & \\
\rowcolor{alignmentred}
(4) allergic purpura (Henoch-Schönlein purpura) & & & \\
\rowcolor{alignmentblue}
B2.2.2 Neoplastic disorders and tumor-like conditions & & & \\
\rowcolor{alignmentyellow}
B2.2.2.1 กลุ่มโรคและกลุ่มอาการที่ต้องรู้ & & & \\
\rowcolor{alignmentred}
B2.2.2.2 กลุ่มโรคและกลุ่มอาการที่ควรรู้ & & & \\
\rowcolor{alignmentred}
(1) lymphoma & & & \\
\rowcolor{alignmentred}
(2) leukemia & & & \\
\rowcolor{alignmentred}
(3) thymoma & & & \\
\rowcolor{alignmentred}
(4) multiple myeloma & & & \\
\rowcolor{alignmentred}
(5) myeloproliferative neoplasms & & & \\
\rowcolor{alignmentblue}
B2.2.3 Metabolic and regulatory disorders & & & \\
\rowcolor{alignmentyellow}
B2.2.3.1 กลุ่มโรคและกลุ่มอาการที่ต้องรู้ & & & \\
\rowcolor{alignmentred}
B2.2.3.2 กลุ่มโรคและกลุ่มอาการที่ควรรู้ & & & \\
\rowcolor{alignmentyellow}
(1) nutritional deficiency anemias (iron, folate, B12) & & & \\
\rowcolor{alignmentred}
(1) acquired hemolytic anemia (paroxysmal nocturnal hematuria) & & & \\
\rowcolor{alignmentyellow}
(2) hereditary hemolytic anemia (thalassemia, G-6-PD deficiency) & & & \\
\rowcolor{alignmentred}
(2) hereditary hemolytic anemia (hemoglobinopathy, hereditary spherocytosis) & & & \\
\rowcolor{alignmentred}
(3) cytopenia (aplastic anemia, agranulocytosis) & & & \\
\rowcolor{alignmentred}
(4) bleeding secondary to coagulation disorder (hemophilia, acquired prothrombin complex deficiency or vitamin K deficiency, consumptive coagulopathy) & & & \\
\rowcolor{alignmentred}
(5) bleeding secondary to platelet disorder (von Willebrand diseases, acquired platelet dysfunction with eosinophilia) & & & \\
\rowcolor{alignmentblue}
B2.2.4 Vascular disorders & & & \\
\rowcolor{alignmentyellow}
B2.2.4.1 กลุ่มโรคและกลุ่มอาการที่ต้องรู้ & & & \\
\rowcolor{alignmentred}
B2.2.4.2 กลุ่มโรคและกลุ่มอาการที่ควรรู้ & & & \\
\rowcolor{alignmentred}
(1) deep vein thrombosis & & & \\
\rowcolor{alignmentblue}
B2.3 มีความรู้เรื่องหลักการรักษาและหลักการใช้ยา โดยเฉพาะยาในบัญชียาหลักแห่งชาติของโรค กลุ่มอาการ & & & \\
\rowcolor{alignmentblue}
และภาวะผิดปกติในข้อ B2.2 และรวมถึงยาในกลุ่มต่อไปนี้ & & & \\
\rowcolor{alignmentblue}
B2.3.1 Whole blood and blood products & & & \\
\rowcolor{alignmentblue}
B2.3.2 Drugs used in nutritional, hypoplastic, hemolytic and renal anemia (anemia of chronic renal failure) & & & \\
\rowcolor{alignmentblue}
B2.3.3 Drugs used in bone marrow transplantation & & & \\
\rowcolor{alignmentblue}
B2.3.4 Anticoagulants, antiplatelet drugs, fibrinolytic drugs and hemostatics & & & \\
\rowcolor{alignmentblue}
B2.4 ตระหนักถึงความส าคัญของพฤติกรรมสุขภาพ เพศ วัย เชื้อชาติ อาชีพ เศรษฐกิจ สังคม วัฒนธรรมและ สิ่งแวดล้อม ตลอดจนปัจจัยอื่นๆ ที่มีผลกระทบต่อการบ าบัดรักษา ฟื้นฟูสภาพ สร้างเสริมสุขภาพและ & & & \\
\rowcolor{alignmentblue}
ป้องกันโรค กลุ่มอาการและภาวะผิดปกติในข้อ B2.2 ซึ่งรวมถึง & & & \\
\rowcolor{alignmentblue}
B2.4.1 Emotional and behavioral factors (e.g., barefoot, diet, depression) & & & \\
\rowcolor{alignmentblue}
B2.4.2 Influence on person, family, and society (e.g., childhood leukemia) & & & \\
\rowcolor{alignmentblue}
B2.4.3 Occupational and other environmental risk factors (e.g., benzene, lead,) & & & \\
\rowcolor{alignmentblue}
B2.4.4 Gender and ethnic factors (hill tribe, herbal treatment) & & & \\
\rowcolor{alignmentblue}
\textbf{B3 Central and Peripheral Nervous Systems} & \textbf{10} & \textbf{30} & TBD \\
\rowcolor{alignmentblue}
B3.1 มีความรู้เรื่องภาวะปกติของระบบประสาท ดังต่อไปนี้ & & & \\
\rowcolor{alignmentblue}
B3.1.1 Embryonic development, fetal maturation, and perinatal changes, including neural tube derivatives, cerebral ventricles, neural crest derivatives & & & \\
\rowcolor{alignmentblue}
B3.1.2 Organ structure and function & & & \\
\rowcolor{alignmentblue}
B3.1.2.1 spinal cord, including gross anatomy, blood supply, and spinal reflexes & & & \\
\rowcolor{alignmentblue}
B3.1.2.2 brain stem, including cranial nerves and nuclei, reticular formation, anatomy, and blood supply & & & \\
\rowcolor{alignmentblue}
B3.1.2.3 brain, including gross anatomy and blood supply; cognition, language, memory; hypothalamic function; limbic system and emotional behavior; circadian rhythms and sleep; control of eye movement & & & \\
\rowcolor{alignmentblue}
B3.1.2.4 sensory systems, including proprioception, pain, itching, vision, hearing, balance, taste, and olfaction & & & \\
\rowcolor{alignmentblue}
B3.1.2.5 motor systems, including brain and spinal cord, basal ganglia and cerebellum & & & \\
\rowcolor{alignmentblue}
B3.1.2. 6 autonomic nervous system & & & \\
\rowcolor{alignmentblue}
B3.1.2.7 peripheral nerve & & & \\
\rowcolor{alignmentblue}
B3.1.3 Cell/tissue structure and function & & & \\
\rowcolor{alignmentblue}
B3.1.3.1 axonal transport & & & \\
\rowcolor{alignmentblue}
B3.1.3.2 excitable properties of neurons, axons and dendrites, including channels synthesis, storage, release, reuptake, and degradation of neurotransmitters and neuromodulators & & & \\
\rowcolor{alignmentblue}
B3.1.3.3 pre- and postsynaptic receptor interactions, trophic and growth factors brain metabolism & & & \\
\rowcolor{alignmentblue}
B3.1.3.4 glia, myelin & & & \\
\rowcolor{alignmentblue}
B3.1.3.5 brain homeostasis; blood-brain barrier, cerebrospinal fluid formation and flow, choroid plexus & & & \\
\rowcolor{alignmentblue}
B3.1.4 Repair, regeneration, and changes associated with stage of life & & & \\
\rowcolor{alignmentblue}
B3.2 มีความรู้เรื่อง สาเหตุ พยาธิก าเนิด พยาธิสรีรวิทยา พยาธิสภาพ ตลอดจนหลักการวินิจฉัยสาเหตุของการ เกิดโรค กลุ่มอาการ หรือภาวะผิดปกติทางระบบประสาทดังต่อไปนี้ & & & \\
\rowcolor{alignmentblue}
B3.2.1 Hereditary, congenital and structural disorders & & & \\
\rowcolor{alignmentyellow}
B3.2.1.1 กลุ่มโรคและกลุ่มอาการที่ต้องรู้ & & & \\
\rowcolor{alignmentred}
B3.2.1.2 กลุ่มโรคและกลุ่มอาการที่ควรรู้ & & & \\
\rowcolor{alignmentred}
(1) encephalocoele, hydrocephalus & & & \\
\rowcolor{alignmentblue}
B3.2.2 Infectious, inflammatory, and immunologic disorders & & & \\
\rowcolor{alignmentyellow}
B3.2.2.1 กลุ่มโรคและกลุ่มอาการที่ต้องรู้ & & & \\
\rowcolor{alignmentred}
B3.2.2.2 กลุ่มโรคและกลุ่มอาการที่ควรรู้ & & & \\
\rowcolor{alignmentyellow}
(1) poliomyelitis & & & \\
\rowcolor{alignmentred}
(1) brain abscess & & & \\
\rowcolor{alignmentyellow}
(2) meningitis & & & \\
\rowcolor{alignmentred}
(2) meningococcemia & & & \\
\rowcolor{alignmentyellow}
(3) encephalitis and myelitis & & & \\
\rowcolor{alignmentred}
(3) myasthenia gravis & & & \\
\rowcolor{alignmentyellow}
(4) rabies & & & \\
\rowcolor{alignmentred}
(4) polyneuropathies (Guillain-Barre syndrome) ) (5) trigeminal neuralgia & & & \\
\rowcolor{alignmentyellow}
(5) tetanus & & & \\
\rowcolor{alignmentyellow}
(6) facial nerve paralysis (Bell’s palsy & & & \\
\rowcolor{alignmentyellow}
(7) acute flaccid paralysis & & & \\
\rowcolor{alignmentyellow}
(8) autoimmune disorders (e.g., Wegener granulomatosis, Goodpasture syndrome) & & & \\
\rowcolor{alignmentblue}
B3.2.3 Traumatic and mechanical disorders & & & \\
\rowcolor{alignmentyellow}
B3.2.3.1 กลุ่มโรคและกลุ่มอาการที่ต้องรู้ & & & \\
\rowcolor{alignmentred}
B3.2.3.2 กลุ่มโรคและกลุ่มอาการที่ควรรู้ & & & \\
\rowcolor{alignmentyellow}
(1) coma & & & \\
\rowcolor{alignmentred}
(1) head and neck injury & & & \\
\rowcolor{alignmentyellow}
(2) convulsion & & & \\
\rowcolor{alignmentred}
(2) spinal cord and peripheral nerve e injury & & & \\
\rowcolor{alignmentyellow}
(3) acute increased intracranial pressur & & & \\
\rowcolor{alignmentyellow}
(4) intracrainal bleeding & & & \\
\rowcolor{alignmentred}
(3) peripheral nerve entrapment (carpal tunnel syndrome, cubital tunnel syndrome) & & & \\
\rowcolor{alignmentred}
(4) hydrocephalus & & & \\
\rowcolor{alignmentblue}
B3.2.4 Neoplastic disorders and tumor-like conditions & & & \\
\rowcolor{alignmentyellow}
B3.2.4.1 กลุ่มโรคและกลุ่มอาการที่ต้องรู้ & & & \\
\rowcolor{alignmentred}
B3.2.4.2 กลุ่มโรคและกลุ่มอาการที่ควรรู้ & & & \\
\rowcolor{alignmentred}
(1) benign and malignant neoplasm of brain & & & \\
\rowcolor{alignmentblue}
B3.2.5 Metabolic and regulatory disorders & & & \\
\rowcolor{alignmentyellow}
B3.2.5.1 กลุ่มโรคและกลุ่มอาการที่ต้องรู้ & & & \\
\rowcolor{alignmentred}
B3.2.5.2 กลุ่มโรคและกลุ่มอาการที่ควรรู้ & & & \\
\rowcolor{alignmentyellow}
(1) periodic paralysis & & & \\
\rowcolor{alignmentred}
(1) Reye’s syndrome & & & \\
\rowcolor{alignmentred}
(2) myopathy & & & \\
\rowcolor{alignmentred}
(3) delirium & & & \\
\rowcolor{alignmentred}
(4) cerebral palsy & & & \\
\rowcolor{alignmentblue}
B3.2.6 Vascular and circulatory disorders & & & \\
\rowcolor{alignmentyellow}
B3.2.6.1 กลุ่มโรคและกลุ่มอาการที่ต้องรู้ & & & \\
\rowcolor{alignmentred}
B3.2.6.2 กลุ่มโรคและกลุ่มอาการที่ควรรู้ & & & \\
\rowcolor{alignmentred}
(1) cerebrovascular diseases (intracerebral hemorrhage, cerebral infarction, subarachnoid hemorrhage) & & & \\
\rowcolor{alignmentblue}
B3.2.7 Degenerative disorders & & & \\
\rowcolor{alignmentyellow}
B3.2.7.1 กลุ่มโรคและกลุ่มอาการที่ต้องรู้ & & & \\
\rowcolor{alignmentred}
B3.2.7.2 กลุ่มโรคและกลุ่มอาการที่ควรรู้ & & & \\
\rowcolor{alignmentred}
(1) Parkinsonism & & & \\
\rowcolor{alignmentred}
(2) Alzheimer’s disease & & & \\
\rowcolor{alignmentred}
(3) polyneuropathies (systemic diseases, drugs, chemical, alcohol) & & & \\
\rowcolor{alignmentred}
(4) peripheral neuropathy & & & \\
\rowcolor{alignmentblue}
B3.2.8 Paroxysmal disorders & & & \\
\rowcolor{alignmentyellow}
B3.2.8.1 กลุ่มโรคและกลุ่มอาการที่ต้องรู้ & & & \\
\rowcolor{alignmentred}
B3.2.8.2 กลุ่มโรคและกลุ่มอาการที่ควรรู้ & & & \\
\rowcolor{alignmentyellow}
(1) epilepsy & & & \\
\rowcolor{alignmentred}
(1) sleep disorders including narcolepsy, restless legs syndrome/periodic limb movement, circadian rhythm disorders, parasomnias & & & \\
\rowcolor{alignmentyellow}
(2) tension headache & & & \\
\rowcolor{alignmentyellow}
(3) migraine headache & & & \\
\rowcolor{alignmentyellow}
(4) pain syndromes & & & \\
\rowcolor{alignmentblue}
B3.2.9 Psychopathologic disorders process and their evaluation & & & \\
\rowcolor{alignmentyellow}
B3.2.9.1 กลุ่มโรคและกลุ่มอาการที่ต้องรู้ & & & \\
\rowcolor{alignmentred}
B3.2.9.2 กลุ่มโรคและกลุ่มอาการที่ควรรู้ & & & \\
\rowcolor{alignmentyellow}
(1) acute psychosis, delirium & & & \\
\rowcolor{alignmentred}
(1) organic mental disorders (dementia, brain damage, systemic diseases, alcohol \& substance induced, organic amnestic syndrome, delirium) & & & \\
\rowcolor{alignmentyellow}
(2) hyperventilation syndrome & & & \\
\rowcolor{alignmentyellow}
(3) severe depression & & & \\
\rowcolor{alignmentyellow}
(4) suicide attempt & & & \\
\rowcolor{alignmentyellow}
(5) panic attack & & & \\
\rowcolor{alignmentyellow}
(6) reaction to severe stress & & & \\
\rowcolor{alignmentred}
(2) neurotic disorders (phobia, obscessive-compulsive o disorders, dissociative or conversion disorder) , (3) somatoform disorders & & & \\
\rowcolor{alignmentyellow}
(7) rape & & & \\
\rowcolor{alignmentyellow}
(8) mental \& behavioral disorders due t alcohol, drugs and substances & & & \\
\rowcolor{alignmentyellow}
(9) depressive disorder, panic disorder generalized anxiety disorders & & & \\
\rowcolor{alignmentred}
(4) behavioral and emotional disorders in childhood and adolescence (conduct disorder, tic disorder, enuresis, feeding disorder, nail biting, thumb sucking, attention deficit disorders) & & & \\
\rowcolor{alignmentyellow}
(10) physical and sexual abuse of children, adults, and elders & & & \\
\rowcolor{alignmentred}
(5) schizophrenia & & & \\
\rowcolor{alignmentred}
(6) mood (affective) disorders; manic episode, bipolar affective disorder & & & \\
\rowcolor{alignmentred}
(7) behavioral syndromes associated with physiological dysfunction and physical factors (eating disorder, sleep disorder, sexual dysfunction, post-partum psychosis) & & & \\
\rowcolor{alignmentred}
(8) disorders of adult personality and behavior (aggressive, antisocial, compulsive, hysterical, obscessive- compulsive, paranoid) & & & \\
\rowcolor{alignmentred}
(9) mental retardation & & & \\
\rowcolor{alignmentred}
(10) disorders of psychological development (speech and language disorder, disorder of scholastic skills, motor function disorder, autism) & & & \\
\rowcolor{alignmentblue}
B3.2.10 Disorder of special senses: Ear & & & \\
\rowcolor{alignmentyellow}
B3.2.10.1 กลุ่มโรคและกลุ่มอาการที่ต้องรู้ & & & \\
\rowcolor{alignmentred}
B3.2.10.2 กลุ่มโรคและกลุ่มอาการที่ควรรู้ & & & \\
\rowcolor{alignmentyellow}
(1) myringitis & & & \\
\rowcolor{alignmentred}
(1) nasal polyps & & & \\
\rowcolor{alignmentyellow}
(2) acute otitis externa & & & \\
\rowcolor{alignmentred}
(2) perforation of tympanic membrane & & & \\
\rowcolor{alignmentyellow}
(3) acute otitis media & & & \\
\rowcolor{alignmentyellow}
(4) impacted cerumen & & & \\
\rowcolor{alignmentred}
(3) chronic otitis media & & & \\
\rowcolor{alignmentyellow}
(5) foreign body through orifice & & & \\
\rowcolor{alignmentred}
(4) mastoiditis, cholesteatoma & & & \\
\rowcolor{alignmentred}
(5) perichondritis of the pinna & & & \\
\rowcolor{alignmentred}
(6) conductive and sensorineural hearing loss & & & \\
\rowcolor{alignmentred}
(7) disorders of vestibular function and vertiginous syndrome (Meniere ‘s syndrome, vertigo) & & & \\
\rowcolor{alignmentred}
(8) benign and malignant neoplasm of larynx, nasopharynx; nasopharyngeal carcinoma, carcinoma of larynx & & & \\
\rowcolor{alignmentblue}
B3.2.11 Disorder of special senses: Eye & & & \\
\rowcolor{alignmentyellow}
B3.2.11.1 กลุ่มโรคและกลุ่มอาการที่ต้องรู้ & & & \\
\rowcolor{alignmentred}
B3.2.11.2 กลุ่มโรคและกลุ่มอาการที่ควรรู้ & & & \\
\rowcolor{alignmentyellow}
(1) eye injury and foreign body on external eye & & & \\
\rowcolor{alignmentred}
(1) disorders of ocular muscles, refraction and accommodation (strabismus, amblyopia, myopia, presbyopia, hypermetropia, astigmatism) & & & \\
\rowcolor{alignmentyellow}
(2) hyphema & & & \\
\rowcolor{alignmentyellow}
(3) acute corneal abrasion and ulcer & & & \\
\rowcolor{alignmentyellow}
(4) acute glaucoma & & & \\
\rowcolor{alignmentyellow}
(5) hordeolum and chalazion & & & \\
\rowcolor{alignmentred}
(2) dacryostenosis, dacryocystitis & & & \\
\rowcolor{alignmentyellow}
(6) conjunctivitis & & & \\
\rowcolor{alignmentred}
(3) pterygium & & & \\
\rowcolor{alignmentyellow}
(7) pinguecula & & & \\
\rowcolor{alignmentred}
(4) keratitis, corneal ulcer & & & \\
\rowcolor{alignmentred}
(5) uveitis & & & \\
\rowcolor{alignmentred}
(6) cataract & & & \\
\rowcolor{alignmentred}
(7) glaucoma & & & \\
\rowcolor{alignmentred}
(8) retinopathy & & & \\
\rowcolor{alignmentred}
(9) benign and malignant neoplasm of eye & & & \\
\rowcolor{alignmentblue}
B3.3 มีความรู้เรื่องหลักการรักษาและหลักการใช้ยา โดยเฉพาะยาในบัญชียาหลักแห่งชาติของโรค กลุ่มอาการ & & & \\
\rowcolor{alignmentblue}
และภาวะผิดปกติในข้อ B3.2 และรวมถึงยาในกลุ่มต่อไปนี้ & & & \\
\rowcolor{alignmentblue}
B3.3.1 Hypnotics and anxiolytics & & & \\
\rowcolor{alignmentblue}
B3.3.2 Drugs used in psychoses and related disorders (antipsychotic drugs and antimanic drugs) & & & \\
\rowcolor{alignmentblue}
B3.3.3 Antidepressant drugs & & & \\
\rowcolor{alignmentblue}
B3.3.4 Central nervous system stimulants & & & \\
\rowcolor{alignmentblue}
B3.3.5 Drugs used in nausea, vertigo and vestibular disorders & & & \\
\rowcolor{alignmentblue}
B3.3.6 Analgesics and antipyretics including opioid analgesics, drugs for neuropathic pain and antimigraine drugs) & & & \\
\rowcolor{alignmentblue}
B3.3.7 Antiepileptics & & & \\
\rowcolor{alignmentblue}
B3.3.8 Drugs used in movement disorders & & & \\
\rowcolor{alignmentblue}
B3.3.9 Drugs used in substance dependence & & & \\
\rowcolor{alignmentblue}
B3.3.10 Drugs used in neuromuscular disorders (drugs enhancing neuromuscular transmission, skeletal muscle relaxants) & & & \\
\rowcolor{alignmentblue}
B3.3.11 Eye preparations & & & \\
\rowcolor{alignmentblue}
B3.3.11.1 anti-infective preparations & & & \\
\rowcolor{alignmentblue}
B3.3.11.2 corticosteroids and other anti-inflammatory preparations & & & \\
\rowcolor{alignmentblue}
B3.3.11.3 mydriatics and cycloplegics & & & \\
\rowcolor{alignmentblue}
B3.3.11.4 drugs for treatment of glaucoma & & & \\
\rowcolor{alignmentblue}
B3.3.11.5 local anesthetics & & & \\
\rowcolor{alignmentblue}
B3.3.12 Drugs acting on the ear & & & \\
\rowcolor{alignmentblue}
B3.4 ตระหนักถึงความส าคัญของพฤติกรรมสุขภาพ เพศ วัย เชื้อชาติ อาชีพ เศรษฐกิจ สังคม วัฒนธรรมและ สิ่งแวดล้อม ตลอดจนปัจจัยอื่นๆ ที่มีผลกระทบต่อการบ าบัดรักษา ฟื้นฟูสภาพ สร้างเสริมสุขภาพและ & & & \\
\rowcolor{alignmentblue}
ป้องกันโรค/กลุ่มอาการ/ภาวะในข้อ B3.2 ซึ่งรวมถึง & & & \\
\rowcolor{alignmentblue}
B3.4.1 Emotional and behavioral factors (e.g., drug abuse, dementia, sleep deprivation, accident prevention, pets) & & & \\
\rowcolor{alignmentblue}
B3.4.2 Influence on person, family, and society (e.g., developmental disabilities, dementia, generation reversal, nutrition, seizures, sleep disorders) & & & \\
\rowcolor{alignmentblue}
B3.4.3 Occupational and other environmental risk factors (e.g., boxing, carbon monoxide exposure) & & & \\
\rowcolor{alignmentblue}
B3.4.4 Gender and ethnic factors & & & \\
\rowcolor{alignmentblue}
\textbf{B4 Skin and Related Connective Tissue} & \textbf{5} & \textbf{15} & TBD \\
\rowcolor{alignmentblue}
B4.1 มีความรู้เรื่องภาวะปกติของผิวหนังและเนื้อเยื่อเกี่ยวพันที่เกี่ยวข้อง ดังต่อไปนี้ & & & \\
\rowcolor{alignmentblue}
B4.1.1 Embryonic development, fetal maturation, and perinatal changes & & & \\
\rowcolor{alignmentblue}
B4.1.2 Organ structure and function & & & \\
\rowcolor{alignmentblue}
B4.1.3 Cell/tissue structure and function, including barrier functions, thermal regulation, eccrine function & & & \\
\rowcolor{alignmentblue}
B4.1.4 Repair, regeneration, and changes associated with stage of life or ethnicity (e.g., senile purpura, male pattern baldness, postmenopausal hair changes) & & & \\
\rowcolor{alignmentblue}
B4.1.5 Skin defense mechanisms and normal flora & & & \\
\rowcolor{alignmentblue}
B4.2. มีความรู้เรื่องสาเหตุ พยาธิก าเนิด พยาธิสรีรวิทยา พยาธิสภาพ ตลอดจนหลักการวินิจฉัยสาเหตุของการ เกิดโรค กลุ่มอาการ หรือภาวะผิดปกติทางผิวหนังและเนื้อเยื่อเกี่ยวพันที่เกี่ยวข้อง ดังต่อไปนี้ & & & \\
\rowcolor{alignmentblue}
B4.2.1 Infectious, inflammatory, and immunologic disorders & & & \\
\rowcolor{alignmentyellow}
B4.2.1.1 กลุ่มโรคและกลุ่มอาการที่ต้องรู้ & & & \\
\rowcolor{alignmentred}
B4.2.1.2 กลุ่มโรคและกลุ่มอาการที่ควรรู้ & & & \\
\rowcolor{alignmentyellow}
(1) acne, impetigo, cellulitis, abscess, & & & \\
\rowcolor{alignmentred}
(1) staphylococcal scalded skin syndrome (SSSS), carbuncle, gangrene & & & \\
\rowcolor{alignmentyellow}
(2) wart (verrucae) & & & \\
\rowcolor{alignmentyellow}
(3) superficial mycoses (tinea, pityriasis versicolor, candidosis) & & & \\
\rowcolor{alignmentred}
(2) leprosy & & & \\
\rowcolor{alignmentyellow}
(4) scabiasis, pediculosis & & & \\
\rowcolor{alignmentred}
(3) bullous dermatoses (pemphigus, pemphigoid) & & & \\
\rowcolor{alignmentyellow}
(5) urticaria, eczema, dermatitis, alopecia & & & \\
\rowcolor{alignmentred}
(4) papulosquamous disorders (psoriasis, pityriasis rosea, lichen planus) & & & \\
\rowcolor{alignmentred}
(5) erythema multiforme, erythema nodosum, Steven-Johnson’s syndrome, toxic epidermal necrolysis & & & \\
\rowcolor{alignmentred}
(6) discoid lupus erythematosus & & & \\
\rowcolor{alignmentred}
(7) scleroderma & & & \\
\rowcolor{alignmentblue}
B4.2.2 Traumatic and mechanical disorders & & & \\
\rowcolor{alignmentyellow}
B4.2.2.1 กลุ่มโรคและกลุ่มอาการที่ต้องรู้ & & & \\
\rowcolor{alignmentred}
B4.2.2.2 กลุ่มโรคและกลุ่มอาการที่ควรรู้ & & & \\
\rowcolor{alignmentyellow}
(1) ulcers & & & \\
\rowcolor{alignmentyellow}
(2) bites, stings & & & \\
\rowcolor{alignmentyellow}
(3) wound & & & \\
\rowcolor{alignmentblue}
B4.2.3 Neoplastic disorders and tumor-like conditions & & & \\
\rowcolor{alignmentyellow}
B4.2.3.1 กลุ่มโรคและกลุ่มอาการที่ต้องรู้ & & & \\
\rowcolor{alignmentred}
B4.2.3.2 กลุ่มโรคและกลุ่มอาการที่ควรรู้ & & & \\
\rowcolor{alignmentyellow}
(1) cyst & & & \\
\rowcolor{alignmentred}
(1) hemangiomas & & & \\
\rowcolor{alignmentyellow}
(2) benign neoplasm of skin and related connective tissue & & & \\
\rowcolor{alignmentred}
(2) malignant neoplasm of skin and related connective tissue & & & \\
\rowcolor{alignmentblue}
B4.2.4 Vascular disorders & & & \\
\rowcolor{alignmentyellow}
B4.2.4.1 กลุ่มโรคและกลุ่มอาการที่ต้องรู้ & & & \\
\rowcolor{alignmentred}
B4.2.4.2 กลุ่มโรคและกลุ่มอาการที่ควรรู้ & & & \\
\rowcolor{alignmentyellow}
(1) vasculitis & & & \\
\rowcolor{alignmentred}
(1) Raynaud’s disease & & & \\
\rowcolor{alignmentblue}
B4.2.5 Others & & & \\
\rowcolor{alignmentyellow}
B4.2.5.1 กลุ่มโรคและกลุ่มอาการที่ต้องรู้ & & & \\
\rowcolor{alignmentred}
B4.2.5.2 กลุ่มโรคและกลุ่มอาการที่ควรรู้ & & & \\
\rowcolor{alignmentyellow}
(1) corn & & & \\
\rowcolor{alignmentyellow}
(2) scar, keloid, & & & \\
\rowcolor{alignmentyellow}
(3) dyshidrosis, miliaria & & & \\
\rowcolor{alignmentyellow}
(4) freckles, melasma, vitiligo & & & \\
\rowcolor{alignmentblue}
B4.3 มีความรู้เรื่องหลักการรักษาและหลักการใช้ยา โดยเฉพาะยาในบัญชียาหลักแห่งชาติของโรค กลุ่มอาการ & & & \\
\rowcolor{alignmentblue}
และภาวะผิดปกติในข้อ B4.2 และรวมถึงยาในกลุ่มต่อไปนี้ & & & \\
\rowcolor{alignmentblue}
B4.3.1 Anti-infective skin preparations & & & \\
\rowcolor{alignmentblue}
B4.3.2 Emollient and barrier preparations & & & \\
\rowcolor{alignmentblue}
B4.3.3 Topical antipruritics & & & \\
\rowcolor{alignmentblue}
B4.3.4 Topical corticosteroids & & & \\
\rowcolor{alignmentblue}
B4.3.5 Other preparations for psoriasis & & & \\
\rowcolor{alignmentblue}
B4.3.6 Preparations for warts and calluses & & & \\
\rowcolor{alignmentblue}
B4.4 ตระหนักถึงความส าคัญของพฤติกรรมสุขภาพ เพศ วัย เชื้อชาติ อาชีพ เศรษฐกิจ สังคม วัฒนธรรมและ สิ่งแวดล้อม ตลอดจนปัจจัยอื่นๆ ที่มีผลกระทบต่อการบ าบัดรักษา ฟื้นฟูสภาพ สร้างเสริมสุขภาพและ & & & \\
\rowcolor{alignmentblue}
ป้องกันโรค/กลุ่มอาการ/ภาวะในข้อ B4.2 รวมถึง & & & \\
\rowcolor{alignmentblue}
B4.4.1 Emotional and behavioral factors (e.g., sun exposure, acne) & & & \\
\rowcolor{alignmentblue}
B4.2.2 Influence on person, family, and society (e.g., psoriasis) & & & \\
\rowcolor{alignmentblue}
B4.4.3 Occupational and other environmental risk factors & & & \\
\rowcolor{alignmentblue}
B4.4.4 Gender and ethnic factors (e.g., keloid) & & & \\
\rowcolor{alignmentblue}
\textbf{B5 Musculoskeletal System} & \textbf{5} & \textbf{15} & TBD \\
\rowcolor{alignmentblue}
B5.1 มีความรู้เรื่องภาวะปกติของระบบกระดูก กล้ามเนื้อ และเนื้อเยื่อเกี่ยวพัน ดังต่อไปนี้ & & & \\
\rowcolor{alignmentblue}
B5.1.1 Embryonic development, fetal maturation, and perinatal changes & & & \\
\rowcolor{alignmentblue}
B5.1.2 Organ structure and function & & & \\
\rowcolor{alignmentblue}
B5.1.3 Cell/tissue structure and function & & & \\
\rowcolor{alignmentblue}
B5.1.3.1 biology of bones, joints, tendons, skeletal muscle & & & \\
\rowcolor{alignmentblue}
B5.1.3.2 exercise and physical conditioning & & & \\
\rowcolor{alignmentblue}
B5.1.4 Repair, regeneration, and changes associated with stage of life & & & \\
\rowcolor{alignmentblue}
B5.2 มีความรู้เรื่อง สาเหตุ พยาธิก าเนิด พยาธิสรีรวิทยา พยาธิสภาพ ตลอดจนหลักการวินิจฉัยสาเหตุของการ เกิดโรค กลุ่มอาการ หรือภาวะผิดปกติทางระบบกระดูก กล้ามเนื้อ และเนื้อเยื่อเกี่ยวพัน ดังต่อไปนี้ & & & \\
\rowcolor{alignmentblue}
B5.2.1 Hereditary, congenital and structural disorders & & & \\
\rowcolor{alignmentyellow}
B5.2.1.1 กลุ่มโรคและกลุ่มอาการที่ต้องรู้ & & & \\
\rowcolor{alignmentred}
B5.2.1.2 กลุ่มโรคและกลุ่มอาการที่ควรรู้ & & & \\
\rowcolor{alignmentred}
(1) clubfoot & & & \\
\rowcolor{alignmentred}
(2) congenital hip dislocation & & & \\
\rowcolor{alignmentblue}
B5.2.2 Infectious, inflammatory, and immunologic disorders & & & \\
\rowcolor{alignmentyellow}
B5.2.2.1 กลุ่มโรคและกลุ่มอาการที่ต้องรู้ & & & \\
\rowcolor{alignmentred}
B5.2.2.2 กลุ่มโรคและกลุ่มอาการที่ควรรู้ & & & \\
\rowcolor{alignmentyellow}
(1) osteomyelitis, & & & \\
\rowcolor{alignmentred}
(1) osteonecrosis & & & \\
\rowcolor{alignmentyellow}
(2) pyomyositis & & & \\
\rowcolor{alignmentred}
(2) spondylitis & & & \\
\rowcolor{alignmentyellow}
(3) soft-tissue rheumatism (e.g., tendinitis, de Quervain diseas bursitis, fasciitis, myofascial pain syndrome) & & & \\
\rowcolor{alignmentred}
(3) rheumatoid arthritis e (4) juvenile arthritis & & & \\
\rowcolor{alignmentred}
(5) spondyloarthropathy & & & \\
\rowcolor{alignmentred}
(6) inflammatory myositis & & & \\
\rowcolor{alignmentyellow}
(4) costochondritis & & & \\
\rowcolor{alignmentred}
(7) necrotizing fasciitis & & & \\
\rowcolor{alignmentyellow}
(5) reactive arthropathy & & & \\
\rowcolor{alignmentyellow}
(6) infective arthritis & & & \\
\rowcolor{alignmentblue}
B5.2.3 Traumatic and mechanical disorders & & & \\
\rowcolor{alignmentyellow}
B5.2.3.1 กลุ่มโรคและกลุ่มอาการที่ต้องรู้ & & & \\
\rowcolor{alignmentred}
B5.2.3.2 กลุ่มโรคและกลุ่มอาการที่ควรรู้ & & & \\
\rowcolor{alignmentyellow}
(1) sprains, strains, fractures, dislocations & & & \\
\rowcolor{alignmentyellow}
(2) compartment syndrome & & & \\
\rowcolor{alignmentblue}
B5.2.4 Neoplastic disorders and tumor-like conditions & & & \\
\rowcolor{alignmentyellow}
B5.2.4.1 กลุ่มโรคและกลุ่มอาการที่ต้องรู้ & & & \\
\rowcolor{alignmentred}
B5.2.4.2 กลุ่มโรคและกลุ่มอาการที่ควรรู้ & & & \\
\rowcolor{alignmentred}
(1) benign and malignant neoplasm of bone and soft tissue & & & \\
\rowcolor{alignmentblue}
B5.2.5 Metabolic and regulatory disorders & & & \\
\rowcolor{alignmentyellow}
B5.2.5.1 กลุ่มโรคและกลุ่มอาการที่ต้องรู้ & & & \\
\rowcolor{alignmentred}
B5.2.5.2 กลุ่มโรคและกลุ่มอาการที่ควรรู้ & & & \\
\rowcolor{alignmentyellow}
(1) crystal arthropathy (gout, pseudogout) & & & \\
\rowcolor{alignmentred}
(2) muscular dystrophy & & & \\
\rowcolor{alignmentblue}
B5.2.6 Vascular disorders & & & \\
\rowcolor{alignmentyellow}
B5.2.6.1 กลุ่มโรคและกลุ่มอาการที่ต้องรู้ & & & \\
\rowcolor{alignmentred}
B5.2.6.2 กลุ่มโรคและกลุ่มอาการที่ควรรู้ & & & \\
\rowcolor{alignmentred}
(1) vasculitis syndrome (Kawasaki’s disease, polyarteritis nodosa, drug induced) & & & \\
\rowcolor{alignmentblue}
B5.2.7 Systemic disorders affecting the musculoskeletal System & & & \\
\rowcolor{alignmentyellow}
B5.2.7.1 กลุ่มโรคและกลุ่มอาการที่ต้องรู้ & & & \\
\rowcolor{alignmentred}
B5.2.7.2 กลุ่มโรคและกลุ่มอาการที่ควรรู้ & & & \\
\rowcolor{alignmentred}
(1) systemic lupus erythematosus & & & \\
\rowcolor{alignmentblue}
B5.2.8 Others & & & \\
\rowcolor{alignmentyellow}
B5.2.8.1 กลุ่มโรคและกลุ่มอาการที่ต้องรู้ & & & \\
\rowcolor{alignmentred}
B5.2.8.2 กลุ่มโรคและกลุ่มอาการที่ควรรู้ & & & \\
\rowcolor{alignmentyellow}
(1) osteoarthritis & & & \\
\rowcolor{alignmentred}
(1) spondylolithiasis and disc syndrome & & & \\
\rowcolor{alignmentyellow}
(2) degenerative diseases of the spine & & & \\
\rowcolor{alignmentred}
(2) osteoporosis & & & \\
\rowcolor{alignmentred}
(3) abnormal curvature of the spine & & & \\
\rowcolor{alignmentred}
(4) fibromyalgia & & & \\
\rowcolor{alignmentblue}
B5.3 มีความรู้เรื่องหลักการรักษาและหลักการใช้ยา โดยเฉพาะยาในบัญชียาหลักแห่งชาติของโรค กลุ่มอาการ & & & \\
\rowcolor{alignmentblue}
และภาวะผิดปกติในข้อ B5.2 และรวมถึงยาในกลุ่มต่อไปนี้ & & & \\
\rowcolor{alignmentblue}
B5.3.1 Non-steroidal anti-inflammatory drugs & & & \\
\rowcolor{alignmentblue}
B5.3.2 Disease-modifying antirheumatic drugs & & & \\
\rowcolor{alignmentblue}
B5.3.3 Drugs for treatment of gout and hyperuricemia & & & \\
\rowcolor{alignmentblue}
B5.3.4 Drugs for relief of soft tissue inflammation (rubifacients) & & & \\
\rowcolor{alignmentblue}
B5.4 ตระหนักถึงความส าคัญของพฤติกรรมสุขภาพ เพศ วัย เชื้อชาติ อาชีพ เศรษฐกิจ สังคม วัฒนธรรมและ สิ่งแวดล้อม ตลอดจนปัจจัยอื่นๆ ที่มีผลกระทบต่อการบ าบัดรักษา ฟื้นฟูสภาพ สร้างเสริมสุขภาพและ & & & \\
\rowcolor{alignmentblue}
ป้องกันโรค/กลุ่มอาการ/ภาวะในข้อ B5.2 ซึ่งรวมถึง & & & \\
\rowcolor{alignmentblue}
B5.4.1 Emotional and behavioral factors (e.g., diet, exercise, seat belts, helmets) & & & \\
\rowcolor{alignmentblue}
B5.4.2 Influence on person, family, and society (e.g., osteoporosis, fractures in elderly, alcohol abuse and fractures) & & & \\
\rowcolor{alignmentblue}
B5.4.3 Occupational and other environmental risk factors (e.g., athletes, musicians) & & & \\
\rowcolor{alignmentblue}
B5.4.4 Gender and ethnic factors (e.g., bone mass) & & & \\
\rowcolor{alignmentblue}
\textbf{B6 Respiratory System} & \textbf{10} & \textbf{30} & TBD \\
\rowcolor{alignmentblue}
B6.1 มีความรู้เรื่องภาวะปกติของระบบหายใจ ดังต่อไปนี้ & & & \\
\rowcolor{alignmentblue}
B6.1.1 Embryonic development, fetal maturation, and perinatal changes & & & \\
\rowcolor{alignmentblue}
B6.1.2 Organ structure and function & & & \\
\rowcolor{alignmentblue}
B6.1.2.1 airways, including mechanics and regulation of breathing & & & \\
\rowcolor{alignmentblue}
B6.1.2.2 lung parenchyma, including ventilation, perfusion, gas exchange & & & \\
\rowcolor{alignmentblue}
B6.1.2.3 pleura & & & \\
\rowcolor{alignmentblue}
B6.1.2.4 nasopharynx and sinuses & & & \\
\rowcolor{alignmentblue}
B6.1.3 Cell/tissue structure and function, including surfactant formation, alveolar structure & & & \\
\rowcolor{alignmentblue}
B6.1.4 Repair, regeneration, and changes associated with stage of life & & & \\
\rowcolor{alignmentblue}
B6.1.5 Pulmonary defense mechanisms and normal flora & & & \\
\rowcolor{alignmentblue}
B6.2 มีความรู้เรื่อง สาเหตุ พยาธิก าเนิด พยาธิสรีรวิทยา พยาธิสภาพ ตลอดจนหลักการวินิจฉัยสาเหตุของการ เกิดโรค กลุ่มอาการ หรือภาวะผิดปกติทางระบบหายใจดังต่อไปนี้ & & & \\
\rowcolor{alignmentblue}
B6.2.1 Hereditary, congenital and structural disorders & & & \\
\rowcolor{alignmentyellow}
B6.2.1.1 กลุ่มโรคและกลุ่มอาการที่ต้องรู้ & & & \\
\rowcolor{alignmentred}
B6.2.1.2 กลุ่มโรคและกลุ่มอาการที่ควรรู้ & & & \\
\rowcolor{alignmentred}
(1) cleft lip and palate & & & \\
\rowcolor{alignmentred}
(2) laryngomalacia & & & \\
\rowcolor{alignmentred}
(3) tracheoesophageal fistula & & & \\
\rowcolor{alignmentred}
(4) diaphragmatic hernia & & & \\
\rowcolor{alignmentred}
(5) pre-auricular sinus, sinus, cyst and fistula of branchial cleft & & & \\
\rowcolor{alignmentred}
(6) thyroglossal duct cyst & & & \\
\rowcolor{alignmentblue}
B6.2.2 Infectious, inflammatory, and immunologic disorders & & & \\
\rowcolor{alignmentyellow}
B6.2.2.1 กลุ่มโรคและกลุ่มอาการที่ต้องรู้ & & & \\
\rowcolor{alignmentred}
B6.2.2.2 กลุ่มโรคและกลุ่มอาการที่ควรรู้ & & & \\
\rowcolor{alignmentyellow}
(1) influenza & & & \\
\rowcolor{alignmentred}
(1) avian flu, SARS & & & \\
\rowcolor{alignmentyellow}
(2) upper respiratory infections (nasopharyngitis, rhinosinusitis, pharyngitis, tonsillitis, laryngitis, tracheitis, croup) & & & \\
\rowcolor{alignmentred}
(2) peritonsillar abscess , deep neck infection & & & \\
\rowcolor{alignmentred}
(3) retropharyngeal abscess & & & \\
\rowcolor{alignmentred}
(4) pyothorax (empyemathoracis) & & & \\
\rowcolor{alignmentyellow}
(3) lower respiratory tract infections and pleura and their complications (bronchitis, bronchiolitis, pneumonia, lung abscess, bronchiectasis) & & & \\
\rowcolor{alignmentred}
(5) chronic infectious diseases of the lower respiratory tract endemic fungal infections, Nocardia/Actinomyces & & & \\
\rowcolor{alignmentyellow}
(4) tuberculosis & & & \\
\rowcolor{alignmentred}
(6) occupational lung diseases & & & \\
\rowcolor{alignmentyellow}
(5) pleurisy & & & \\
\rowcolor{alignmentred}
(7) acute and chronic alveolar injury (e.g., acute respiratory distress syndrome, chlorine gas/smoke inhalation) & & & \\
\rowcolor{alignmentyellow}
(6) allergic rhinitis & & & \\
\rowcolor{alignmentyellow}
(7) asthma & & & \\
\rowcolor{alignmentyellow}
(8) autoimmune disorders (e.g., Wegener granulomatosis, Goodpasture syndrome) & & & \\
\rowcolor{alignmentred}
(8) obstructive pulmonary disease (chronic bronchitis, emphysema) & & & \\
\rowcolor{alignmentred}
(9) restrictive pulmonary disease & & & \\
\rowcolor{alignmentblue}
B6.2.3 Traumatic and mechanical disorders & & & \\
\rowcolor{alignmentyellow}
B6.2.3.1 กลุ่มโรคและกลุ่มอาการที่ต้องรู้ & & & \\
\rowcolor{alignmentred}
B6.2.3.2 กลุ่มโรคและกลุ่มอาการที่ควรรู้ & & & \\
\rowcolor{alignmentyellow}
(1) chest injury & & & \\
\rowcolor{alignmentred}
(1) sleep-disordered breathing & & & \\
\rowcolor{alignmentyellow}
(2) foreign body aspiration & & & \\
\rowcolor{alignmentyellow}
(3) pneumothorax & & & \\
\rowcolor{alignmentyellow}
(4) hemothorax & & & \\
\rowcolor{alignmentyellow}
(5) atelectasis & & & \\
\rowcolor{alignmentblue}
B6.2.4 Neoplastic disorders and tumor-like conditions & & & \\
\rowcolor{alignmentyellow}
B6.2.4.1 กลุ่มโรคและกลุ่มอาการที่ต้องรู้ & & & \\
\rowcolor{alignmentred}
B6.2.4.2 กลุ่มโรคและกลุ่มอาการที่ควรรู้ & & & \\
\rowcolor{alignmentred}
(1) benign and malignant neoplasm of upper and lower respiratory tracts & & & \\
\rowcolor{alignmentblue}
B6.2.5 Metabolic, regulatory disorders & & & \\
\rowcolor{alignmentyellow}
B6.2.5.1 กลุ่มโรคและกลุ่มอาการที่ต้องรู้ & & & \\
\rowcolor{alignmentred}
B6.2.5.2 กลุ่มโรคและกลุ่มอาการที่ควรรู้ & & & \\
\rowcolor{alignmentyellow}
(1) respiratory failure & & & \\
\rowcolor{alignmentred}
(1) neonatal respiratory distress syndrome & & & \\
\rowcolor{alignmentyellow}
(2) hyperventilation syndrome & & & \\
\rowcolor{alignmentyellow}
(3) respiratory obstruction, suffocation & & & \\
\rowcolor{alignmentred}
(2) bronchopulmonary dysplasia & & & \\
\rowcolor{alignmentyellow}
(4) asphyxia of the newborn & & & \\
\rowcolor{alignmentyellow}
(5) fetal distress & & & \\
\rowcolor{alignmentyellow}
(6) hypoventilation, disorders of gas exchange, ventilation-perfusion imbalance & & & \\
\rowcolor{alignmentblue}
B6.2.6 Vascular and circulatory disorders & & & \\
\rowcolor{alignmentyellow}
B6.2.6.1 กลุ่มโรคและกลุ่มอาการที่ต้องรู้ & & & \\
\rowcolor{alignmentred}
B6.2.6.2 กลุ่มโรคและกลุ่มอาการที่ควรรู้ & & & \\
\rowcolor{alignmentyellow}
(1) pulmonary edema, pleural effusion & & & \\
\rowcolor{alignmentred}
(1) pulmonary embolism & & & \\
\rowcolor{alignmentred}
(2) pulmonary hypertension & & & \\
\rowcolor{alignmentblue}
B6.3 มีความรู้เรื่องหลักการรักษาและหลักการใช้ยา โดยเฉพาะยาในบัญชียาหลักแห่งชาติของโรค กลุ่มอาการ & & & \\
\rowcolor{alignmentblue}
และภาวะผิดปกติในข้อ B6.2 และรวมถึงยาในกลุ่มต่อไปนี้ & & & \\
\rowcolor{alignmentblue}
B6.3.1 Bronchodilators (adrenoceptor agonists, compound antimuscarinic bronchodilators and theophylline) & & & \\
\rowcolor{alignmentblue}
B6.3.2 Inhaled corticosteroids & & & \\
\rowcolor{alignmentblue}
B6.3.3 Leukotriene receptor antagonists & & & \\
\rowcolor{alignmentblue}
B6.3.4 Antihistamines & & & \\
\rowcolor{alignmentblue}
B6.3.5 Pulmonary surfactants & & & \\
\rowcolor{alignmentblue}
B6.3.6 Cough preparations (cough suppressants, mucolytics, expectorant and demulcent cough preparations) & & & \\
\rowcolor{alignmentblue}
B6.3.7 Systemic and topical nasal decongestants & & & \\
\rowcolor{alignmentblue}
B6.3.8 Other respiratory preparations (aromatic ammonia spirit) & & & \\
\rowcolor{alignmentblue}
B6.3.10 Drugs used in nasal allergy & & & \\
\rowcolor{alignmentblue}
B6.3.11 Antituberculous drugs & & & \\
\rowcolor{alignmentblue}
B6.4 ตระหนักถึงความส าคัญของพฤติกรรมสุขภาพ เพศ วัย เชื้อชาติ อาชีพ เศรษฐกิจ สังคม วัฒนธรรมและ สิ่งแวดล้อม ตลอดจนปัจจัยอื่นๆ ที่มีผลกระทบต่อการบ าบัดรักษา ฟื้นฟูสภาพ สร้างเสริมสุขภาพและ & & & \\
\rowcolor{alignmentblue}
ป้องกันโรค/กลุ่มอาการ/ภาวะในข้อ B6.2 ซึ่งรวมถึง & & & \\
\rowcolor{alignmentblue}
B6.4.1 Emotional and behavioral factors (e.g., smoking, substance abuse, pets, and allergies) & & & \\
\rowcolor{alignmentblue}
B6.4.2 Iobstructive pulmonarynfluence on person, family, and society (e.g., tuberculosis, asthma, chronic disease, school issues, protective parents, family smoking) & & & \\
\rowcolor{alignmentblue}
B6.4.3 Occupational and other environmental risk factors (occupational lung diseases) & & & \\
\rowcolor{alignmentblue}
B6.4.4 Gender and ethnic factors (e.g., lung cancer) & & & \\
\rowcolor{alignmentblue}
\textbf{B7 Cardiovascular System} & \textbf{10} & \textbf{30} & TBD \\
\rowcolor{alignmentblue}
B7.1 มีความรู้เรื่องภาวะปกติของระบบหัวใจหลอดเลือด ดังต่อไปนี้ & & & \\
\rowcolor{alignmentblue}
B7.1.1 Embryonic development, fetal maturation, and perinatal changes & & & \\
\rowcolor{alignmentblue}
B7.1.2 Organ structure and function & & & \\
\rowcolor{alignmentblue}
B7.1.2.1 chambers, valves & & & \\
\rowcolor{alignmentblue}
B7.1.2.2 cardiac cycle, mechanics, heart sounds, cardiac conduction & & & \\
\rowcolor{alignmentblue}
B7.1.2.3 hemodynamics, including systemic, pulmonary, coronary, and blood volume & & & \\
\rowcolor{alignmentblue}
B7.1.2.4 circulation in specific vascular beds & & & \\
\rowcolor{alignmentblue}
B7.1.3 Cell/tissue structure and function & & & \\
\rowcolor{alignmentblue}
B7.1.3.1 heart muscle, metabolism, oxygen consumption, biochemistry, and secretory function (e.g., atrial natriuretic peptide) & & & \\
\rowcolor{alignmentblue}
B7.1.3.2 endothelium and secretory function, vascular smooth muscle, microcirculation, and lymph flow & & & \\
\rowcolor{alignmentblue}
B7.1.3.3 mechanisms of atherosclerosis & & & \\
\rowcolor{alignmentblue}
B7.1.3.4 neural and hormonal regulation of the heart, blood vessels, and blood volume, including responses to change in posture, exercise, and tissue metabolism & & & \\
\rowcolor{alignmentblue}
B7.1.4 Repair, regeneration, and changes associated with stage of life & & & \\
\rowcolor{alignmentblue}
B7.2 มีความรู้เรื่องสาเหตุ พยาธิก าเนิด พยาธิสรีรวิทยา พยาธิสภาพ ตลอดจนหลักการวินิจฉัยสาเหตุของการเกิดโรคกลุ่ม อาการและภาวะผิดปกติทางระบบหัวใจหลอดเลือด ดังต่อไปนี้ & & & \\
\rowcolor{alignmentblue}
B7.2.1 Hereditary, Congenital and structural disorders & & & \\
\rowcolor{alignmentyellow}
B7.2.1.1 กลุ่มโรคและกลุ่มอาการที่ต้องรู้ & & & \\
\rowcolor{alignmentred}
B7.2.1.2 กลุ่มโรคและกลุ่มอาการที่ควรรู้ & & & \\
\rowcolor{alignmentred}
(1) congenital malformation of heart & & & \\
\rowcolor{alignmentblue}
B7.2.2 Infectious, inflammatory, and immunologic disorders & & & \\
\rowcolor{alignmentyellow}
B7.2.2.1 กลุ่มโรคและกลุ่มอาการที่ต้องรู้ & & & \\
\rowcolor{alignmentred}
B7.2.2.2 กลุ่มโรคและกลุ่มอาการที่ควรรู้ & & & \\
\rowcolor{alignmentyellow}
(1) acute rheumatic fever & & & \\
\rowcolor{alignmentred}
(1) infective endocarditis & & & \\
\rowcolor{alignmentred}
(2) myocarditis & & & \\
\rowcolor{alignmentred}
(3) pericarditis & & & \\
\rowcolor{alignmentblue}
B7.2.3 Traumatic and mechanical disorders & & & \\
\rowcolor{alignmentyellow}
B7.2.3.1 กลุ่มโรคและกลุ่มอาการที่ต้องรู้ & & & \\
\rowcolor{alignmentred}
B7.2.3.2 กลุ่มโรคและกลุ่มอาการที่ควรรู้ & & & \\
\rowcolor{alignmentyellow}
(1) anoxic spell & & & \\
\rowcolor{alignmentred}
(1) valvular heart diseases (mitral, aortic valve) & & & \\
\rowcolor{alignmentyellow}
(2) cardiac tamponade & & & \\
\rowcolor{alignmentyellow}
(3) superior vena cava obstruction & & & \\
\rowcolor{alignmentblue}
B7.2.4 Metabolic and regulatory disorders & & & \\
\rowcolor{alignmentyellow}
B7.2.4.1 กลุ่มโรคและกลุ่มอาการที่ต้องรู้ & & & \\
\rowcolor{alignmentred}
B7.2.4.2 กลุ่มโรคและกลุ่มอาการที่ควรรู้ & & & \\
\rowcolor{alignmentyellow}
(1) hypertension & & & \\
\rowcolor{alignmentred}
(1) hypertensive disorder in pregnancy & & & \\
\rowcolor{alignmentyellow}
(2) malignant hypertension & & & \\
\rowcolor{alignmentred}
(2) ischemic heart diseases (angina pectoris, myocardial infarction) & & & \\
\rowcolor{alignmentyellow}
(3) acute coronary syndrome & & & \\
\rowcolor{alignmentyellow}
(4) heart failure & & & \\
\rowcolor{alignmentred}
(3) common cardiac arrhythmias - atrial fibrillation - supraventricular tachycardia - premature ventricular contraction - premature atrial contraction - ventricular tachycardia - ventricular fibrillation - heart block & & & \\
\rowcolor{alignmentyellow}
(5) syncope & & & \\
\rowcolor{alignmentyellow}
(6) cardiac arrest & & & \\
\rowcolor{alignmentyellow}
(7) shock : hypovolemia, cardiogenic & & & \\
\rowcolor{alignmentblue}
B7.2.5 Vascular disorders & & & \\
\rowcolor{alignmentyellow}
B7.2.5.1 กลุ่มโรคและกลุ่มอาการที่ต้องรู้ & & & \\
\rowcolor{alignmentred}
B7.2.5.2 กลุ่มโรคและกลุ่มอาการที่ควรรู้ & & & \\
\rowcolor{alignmentred}
(1) aortic aneurysm and dissection & & & \\
\rowcolor{alignmentred}
(2) varicose veins of lower extremity & & & \\
\rowcolor{alignmentred}
(3) peripheral vascular occlusive disease & & & \\
\rowcolor{alignmentblue}
B7.3 มีความรู้เรื่องหลักการรักษาและหลักการใช้ยา โดยเฉพาะยาในบัญชียาหลักแห่งชาติของโรค กลุ่มอาการ & & & \\
\rowcolor{alignmentblue}
และภาวะผิดปกติในข้อ B7.2 และรวมถึงยาในกลุ่มต่อไปนี้ & & & \\
\rowcolor{alignmentblue}
B7.3.1 Positive inotropic drugs & & & \\
\rowcolor{alignmentblue}
B7.3.2 Diuretics & & & \\
\rowcolor{alignmentblue}
B7.3.3 Anti-arrhythmic drugs & & & \\
\rowcolor{alignmentblue}
B7.3.4 Beta-adrenoceptor blocking drugs & & & \\
\rowcolor{alignmentblue}
B7.3.5 Drugs affecting the renin-angiotensin system and some other antihypertensive drugs (vasodilator antihypertensive drugs, centrally acting antihypertensive drugs, alpha- adrenoceptor blocking drugs, angiotensin-converting enzyme inhibitors, angiotensin-II receptor antagonists) & & & \\
\rowcolor{alignmentblue}
B7.3.6 Nitrates, calcium-channel blockers and other vasodilators & & & \\
\rowcolor{alignmentblue}
B7.3.7 Inotropic sympathomimetics, vasoconstrictor sympathomimetics and drugs used in cardiopulmonary resuscitation & & & \\
\rowcolor{alignmentblue}
B7.4 ตระหนักถึงความส าคัญของพฤติกรรมสุขภาพ เพศ วัย เชื้อชาติ อาชีพ เศรษฐกิจ สังคม วัฒนธรรมและ สิ่งแวดล้อม ตลอดจนปัจจัยอื่นๆ ที่มีผลกระทบต่อการบ าบัดรักษา ฟื้นฟูสภาพ สร้างเสริมสุขภาพและ & & & \\
\rowcolor{alignmentblue}
ป้องกันโรค กลุ่มอาการและภาวะผิดปกติในข้อ B7.2 ซึ่งรวมถึง & & & \\
\rowcolor{alignmentblue}
B7.4.1 Emotional and behavioral factors (e.g., diet, obesity, exercise, smoking, alcohol, stress) & & & \\
\rowcolor{alignmentblue}
B7.4.2 Influence on person, family and society (e.g., lifestyle modification) & & & \\
\rowcolor{alignmentblue}
B7.4.3 Occupational and other environmental risk factors (e.g., stress) & & & \\
\rowcolor{alignmentblue}
B7.4.4 Gender and ethnic factors (e.g., hypertension) & & & \\
\rowcolor{alignmentblue}
\textbf{B8 Gastrointestinal System} & \textbf{10} & \textbf{30} & TBD \\
\rowcolor{alignmentblue}
B8.1 มีความรู้เรื่องภาวะปกติของระบบทางเดินอาหาร ดังต่อไปนี้ & & & \\
\rowcolor{alignmentblue}
B8.1.1 Embryonic development, fetal maturation, and perinatal changes & & & \\
\rowcolor{alignmentblue}
B8.1.2 Organ structure and function function, including alimentary canal, liver and biliary system, salivary glands and exocrine pancreas, motility, and digestion and absorption & & & \\
\rowcolor{alignmentblue}
B8.1.3 Cell/tissue structure and function & & & \\
\rowcolor{alignmentblue}
B8.1.3.1 endocrine and neural regulatory functions, including GI hormones & & & \\
\rowcolor{alignmentblue}
B8.1.3.2 salivary, gastrointestinal, pancreatic, hepatic secretory products, including enzymes, proteins, bile salts, and processes & & & \\
\rowcolor{alignmentblue}
B8.1.3.3 synthetic and metabolic functions of hepatocytes & & & \\
\rowcolor{alignmentblue}
B8.1.4 Repair, regeneration, and changes associated with stage of life & & & \\
\rowcolor{alignmentblue}
B8.1.5 Gastrointestinal defense mechanisms and normal flora & & & \\
\rowcolor{alignmentblue}
B8.2 มีความรู้เรื่องสาเหตุ พยาธิก าเนิด พยาธิสรีรวิทยา พยาธิสภาพ ตลอดจนหลักการวินิจฉัยสาเหตุของการ เกิดโรค กลุ่มอาการ และภาวะผิดปกติทางระบบทางเดินอาหาร ดังต่อไปนี้ & & & \\
\rowcolor{alignmentblue}
B8.2.1 Hereditary, congenital and structural disorders & & & \\
\rowcolor{alignmentyellow}
B8.2.1.1 กลุ่มโรคและกลุ่มอาการที่ต้องรู้ & & & \\
\rowcolor{alignmentred}
B8.2.1.2 กลุ่มโรคและกลุ่มอาการที่ควรรู้ & & & \\
\rowcolor{alignmentred}
(1) disorders of tooth development \& eruption & & & \\
\rowcolor{alignmentred}
(2) hernia (inguinal, umbilical) & & & \\
\rowcolor{alignmentred}
(4) tracheo-esophageal fistula & & & \\
\rowcolor{alignmentred}
(5) diaphragmatic hernia & & & \\
\rowcolor{alignmentred}
(6) Hirschprung diseases, imperforate anus, volvulus & & & \\
\rowcolor{alignmentred}
(7) omphalocoele, gastroschisis & & & \\
\rowcolor{alignmentblue}
B8.2.2 Infectious, inflammatory, and immunologic disorders & & & \\
\rowcolor{alignmentyellow}
B8.2.2.1 กลุ่มโรคและกลุ่มอาการที่ต้องรู้ & & & \\
\rowcolor{alignmentred}
B8.2.2.2 กลุ่มโรคและกลุ่มอาการที่ควรรู้ & & & \\
\rowcolor{alignmentyellow}
(1) gastroenteritis (cholera, shigellosis, salmonellosis, E.coli, Campylobacter sp, Clostridium difficile, virus, amoebiasis , giardiasis, cryptosporidiosis, isosporosis, microsporidiosis, food-borne intoxication) & & & \\
\rowcolor{alignmentred}
(1) dental caries, gingivitis and periodontal diseases & & & \\
\rowcolor{alignmentred}
(2) gastroenteritis in compromised host (cryptosporidiosis, isosporosis, microsporidiosis) & & & \\
\rowcolor{alignmentred}
(3) chronic diarrhea & & & \\
\rowcolor{alignmentred}
(4) necrotizing enterocolitis & & & \\
\rowcolor{alignmentyellow}
(2) typhoid and paratyphoid fever & & & \\
\rowcolor{alignmentred}
(5) peritonitis & & & \\
\rowcolor{alignmentyellow}
(3) intestinal parasites (ascariasis, hookworm infestation, enterobiasis, trichuriasis, strongyloidiasis, taenias capillariasis) & & & \\
\rowcolor{alignmentred}
(6) chronic viral hepatitis & & & \\
\rowcolor{alignmentred}
(7) cholecystitis, cholangitis is, (8) perianal abscess & & & \\
\rowcolor{alignmentyellow}
(4) acute appendicitis & & & \\
\rowcolor{alignmentyellow}
(5) oral ulcer (aphthous ulcer, stomatitis, glossitis, thrush) & & & \\
\rowcolor{alignmentyellow}
(6) peptic ulcer, gastritis & & & \\
\rowcolor{alignmentyellow}
(7) dyspepsia, esophagitis & & & \\
\rowcolor{alignmentyellow}
(8) acute pancreatitis & & & \\
\rowcolor{alignmentyellow}
(9) acute viral hepatitis & & & \\
\rowcolor{alignmentyellow}
(10) alcoholic liver disease & & & \\
\rowcolor{alignmentyellow}
(11) liver abscess (12 ) cirrhosis & & & \\
\rowcolor{alignmentblue}
B8.2.3 Traumatic and mechanical disorders & & & \\
\rowcolor{alignmentyellow}
B8.2.3.1 กลุ่มโรคและกลุ่มอาการที่ต้องรู้ & & & \\
\rowcolor{alignmentred}
B8.2.3.2 กลุ่มโรคและกลุ่มอาการที่ควรรู้ & & & \\
\rowcolor{alignmentyellow}
(1) anal fissure & & & \\
\rowcolor{alignmentred}
(1) intestinal obstruction & & & \\
\rowcolor{alignmentyellow}
(2) gastro-esophageal reflux & & & \\
\rowcolor{alignmentred}
(2) abdominal injury & & & \\
\rowcolor{alignmentred}
(3) infantile hypertrophic pyloric stenosis & & & \\
\rowcolor{alignmentred}
(4) fistula in ano & & & \\
\rowcolor{alignmentblue}
B8.2.4 Neoplastic disorders and tumor-like conditions & & & \\
\rowcolor{alignmentyellow}
B8.2.4.1 กลุ่มโรคและกลุ่มอาการที่ต้องรู้ & & & \\
\rowcolor{alignmentred}
B8.2.4.2 กลุ่มโรคและกลุ่มอาการที่ควรรู้ & & & \\
\rowcolor{alignmentred}
(1) benign and malignant neoplasm of oral cavity, esophagus, stomach, colon, liver and biliary tract, pancreas & & & \\
\rowcolor{alignmentblue}
B8.2.5 Metabolic and regulatory disorders & & & \\
\rowcolor{alignmentyellow}
B8.2.5.1 กลุ่มโรคและกลุ่มอาการที่ต้องรู้ & & & \\
\rowcolor{alignmentred}
B8.2.5.2 กลุ่มโรคและกลุ่มอาการที่ควรรู้ & & & \\
\rowcolor{alignmentyellow}
(1) irritable bowel syndrome & & & \\
\rowcolor{alignmentred}
(1) cholelithiasis & & & \\
\rowcolor{alignmentred}
(2) hepatic failure & & & \\
\rowcolor{alignmentblue}
B8.2.6 Vascular disorders & & & \\
\rowcolor{alignmentyellow}
B8.2.6.1 กลุ่มโรคและกลุ่มอาการที่ต้องรู้ & & & \\
\rowcolor{alignmentred}
B8.2.6.2 กลุ่มโรคและกลุ่มอาการที่ควรรู้ & & & \\
\rowcolor{alignmentyellow}
(1) hemorrhoids & & & \\
\rowcolor{alignmentyellow}
(2) gastrointestinal hemorrhage & & & \\
\rowcolor{alignmentblue}
B8.3 มีความรู้เรื่องหลักการรักษาและหลักการใช้ยา โดยเฉพาะยาในบัญชียาหลักแห่งชาติของโรค กลุ่มอาการและ & & & \\
\rowcolor{alignmentblue}
ภาวะผิดปกติในข้อ B8.2 และรวมถึงยาในกลุ่มต่อไปนี้ & & & \\
\rowcolor{alignmentblue}
B8.3.1 Antacids and other drugs for dyspepsia & & & \\
\rowcolor{alignmentblue}
B8.3.2 Antispasmodics and other drugs altering gut motility & & & \\
\rowcolor{alignmentblue}
B8.3.3 Ulcer-healing drugs and drugs used in variceal bleeding & & & \\
\rowcolor{alignmentblue}
B8.3.4 Drugs used in acute diarrhea & & & \\
\rowcolor{alignmentblue}
B8.3.5 Drugs used in chronic bowel disorders & & & \\
\rowcolor{alignmentblue}
B8.3.6 Laxatives & & & \\
\rowcolor{alignmentblue}
B8.3.7 Local preparations for anal and rectal disorders & & & \\
\rowcolor{alignmentblue}
B8.3.8 Drugs affecting intestinal secretions & & & \\
\rowcolor{alignmentblue}
B8.3.9 Drugs used in treatment of oral ulcer & & & \\
\rowcolor{alignmentblue}
B8.4 ตระหนักถึงความส าคัญของพฤติกรรมสุขภาพ เพศ วัย เชื้อชาติ อาชีพ เศรษฐกิจ สังคม วัฒนธรรมและ สิ่งแวดล้อม ตลอดจนปัจจัยอื่นๆ ที่มีผลกระทบต่อการบ าบัดรักษา ฟื้นฟูสภาพ สร้างเสริมสุขภาพและป้องกัน & & & \\
\rowcolor{alignmentblue}
โรค กลุ่มอาการและภาวะผิดปกติในข้อ B8.2 ซึ่งรวมถึง & & & \\
\rowcolor{alignmentblue}
B8.4.1 Emotional and behavioral factors (e.g., peptic ulcer, encopresis, Monday morning stomach) & & & \\
\rowcolor{alignmentblue}
B8.4.2 Influence on person, family, and society (e.g., inflammatory bowel disease, irritable bowel disease, pancreatitis and alcohol, chronic laxative abuse) & & & \\
\rowcolor{alignmentblue}
B8.4.3 Occupational and other environmental risk factors & & & \\
\rowcolor{alignmentblue}
B8.4.4 Gender and ethnic factors (e.g., diets) & & & \\
\rowcolor{alignmentblue}
\textbf{B9 Renal/Urinary System} & \textbf{10} & \textbf{30} & TBD \\
\rowcolor{alignmentblue}
B9.1 มีความรู้เรื่องภาวะปกติของไตและระบบปัสสาวะดังต่อไปนี้ & & & \\
\rowcolor{alignmentblue}
B9.1.1 Embryonic development, fetal maturation, and perinatal changes & & & \\
\rowcolor{alignmentblue}
B9.1.2 Organ structure and function & & & \\
\rowcolor{alignmentblue}
B9.1.2.1 kidneys, ureters, bladder, urethra & & & \\
\rowcolor{alignmentblue}
B9.1.2.2 glomerular filtration and hemodynamics & & & \\
\rowcolor{alignmentblue}
B9.1.2.3 tubular reabsorption and secretion, including transport processes and proteins & & & \\
\rowcolor{alignmentblue}
B9.1.2.4 urinary concentration and dilution & & & \\
\rowcolor{alignmentblue}
B9.1.2.5 renal mechanisms in acid-base balance & & & \\
\rowcolor{alignmentblue}
B9.1.2.6 renal mechanisms in body fluid homeostasis & & & \\
\rowcolor{alignmentblue}
B9.1.2.7 micturition & & & \\
\rowcolor{alignmentblue}
B9.1.3 Cell/tissue structure and function including renal metabolism and oxygen consumption, hormones produced by or acting on the kidney & & & \\
\rowcolor{alignmentblue}
B9.1.4 Repair, regeneration, and changes associated with stage of life & & & \\
\rowcolor{alignmentblue}
B9.2 มีความรู้เรื่องสาเหตุ พยาธิก าเนิด พยาธิสรีรวิทยา พยาธิสภาพ ตลอดจนหลักการวินิจฉัยหาสาเหตุของการ เกิดโรค กลุ่มอาการและภาวะผิดปกติทางไตและระบบปัสสาวะดังต่อไปนี้ & & & \\
\rowcolor{alignmentblue}
B9.2.1 Hereditary, congenital and structural disorders & & & \\
\rowcolor{alignmentyellow}
B9.2.1.1 กลุ่มโรคและกลุ่มอาการที่ต้องรู้ & & & \\
\rowcolor{alignmentred}
B9.2.1.2 กลุ่มโรคและกลุ่มอาการที่ควรรู้ & & & \\
\rowcolor{alignmentred}
(1) polycystic kidney & & & \\
\rowcolor{alignmentblue}
B9.2.2 Infectious, inflammatory, and immunologic disorders & & & \\
\rowcolor{alignmentyellow}
B9.2.2.1 กลุ่มโรคและกลุ่มอาการที่ต้องรู้ & & & \\
\rowcolor{alignmentred}
B9.2.2.2 กลุ่มโรคและกลุ่มอาการที่ควรรู้ & & & \\
\rowcolor{alignmentyellow}
(1) pyelonephritis & & & \\
\rowcolor{alignmentred}
(1) chronic glomerulonephritis & & & \\
\rowcolor{alignmentyellow}
(2) cystitis, urtethritis & & & \\
\rowcolor{alignmentred}
(2) interstitial nephritis & & & \\
\rowcolor{alignmentyellow}
(3) acute glomerulonephritis & & & \\
\rowcolor{alignmentyellow}
(4) nephrotic syndrome & & & \\
\rowcolor{alignmentblue}
B9.2.3 Traumatic and mechanical disorders & & & \\
\rowcolor{alignmentyellow}
B9.2.3.1 กลุ่มโรคและกลุ่มอาการที่ต้องรู้ & & & \\
\rowcolor{alignmentred}
B9.2.3.2 กลุ่มโรคและกลุ่มอาการที่ควรรู้ & & & \\
\rowcolor{alignmentyellow}
(1) acute urinary retention & & & \\
\rowcolor{alignmentred}
(1) obstructive \& reflux uropathy & & & \\
\rowcolor{alignmentyellow}
(2) urethral syndrome (e.g., urethral stricture) & & & \\
\rowcolor{alignmentred}
(2) urolithiasis & & & \\
\rowcolor{alignmentred}
(3) genitourinary tract injury & & & \\
\rowcolor{alignmentblue}
B9.2.4 Neoplastic disorders and tumor-like condition & & & \\
\rowcolor{alignmentyellow}
B9.2.4.1 กลุ่มโรคและกลุ่มอาการที่ต้องรู้ & & & \\
\rowcolor{alignmentred}
B9.2.4.2 กลุ่มโรคและกลุ่มอาการที่ควรรู้ & & & \\
\rowcolor{alignmentred}
(1) benign and malignant neoplasm of kidney and urinary bladder & & & \\
\rowcolor{alignmentblue}
B9.2.5 Metabolic and regulatory disorders & & & \\
\rowcolor{alignmentyellow}
B9.2.5.1 กลุ่มโรคและกลุ่มอาการที่ต้องรู้ & & & \\
\rowcolor{alignmentred}
B9.2.5.2 กลุ่มโรคและกลุ่มอาการที่ควรรู้ & & & \\
\rowcolor{alignmentyellow}
(1) anuria/oliguria & & & \\
\rowcolor{alignmentred}
(1) renal tubular acidosis & & & \\
\rowcolor{alignmentyellow}
(2) hyperkalemia & & & \\
\rowcolor{alignmentred}
(2) renal failure & & & \\
\rowcolor{alignmentblue}
B9.2.6 Vascular disorders & & & \\
\rowcolor{alignmentyellow}
B9.2.6.1 กลุ่มโรคและกลุ่มอาการที่ต้องรู้ & & & \\
\rowcolor{alignmentred}
B9.2.6.2 กลุ่มโรคและกลุ่มอาการที่ควรรู้ & & & \\
\rowcolor{alignmentred}
(1) renal artery stenosis & & & \\
\rowcolor{alignmentblue}
B9.2.7 Systemic disorders affecting the Renal/Urinary system & & & \\
\rowcolor{alignmentyellow}
B9.2.7.1 กลุ่มโรคและกลุ่มอาการที่ต้องรู้ & & & \\
\rowcolor{alignmentred}
B9.2.7.2 กลุ่มโรคและกลุ่มอาการที่ควรรู้ & & & \\
\rowcolor{alignmentred}
(1) diabetic nephropathy & & & \\
\rowcolor{alignmentred}
(2) lupus nephritis & & & \\
\rowcolor{alignmentblue}
B9.3 มีความรู้เรื่องหลักการรักษาและหลักการใช้ยา โดยเฉพาะยาในบัญชียาหลักแห่งชาติของโรค กลุ่มอาการและ & & & \\
\rowcolor{alignmentblue}
ภาวะผิดปกติในข้อ B9.2 และรวมถึงยาในกลุ่มต่อไปนี้ & & & \\
\rowcolor{alignmentblue}
B9.3.1 Diuretics & & & \\
\rowcolor{alignmentblue}
B9.3.2 Drugs and fluids used to treat volume, electrolyte, and acid-base disorders & & & \\
\rowcolor{alignmentblue}
B9.3.3 Drugs for benign prostatic hyperplasia & & & \\
\rowcolor{alignmentblue}
B9.3.4 Drugs for urinary frequency, enuresis and incontinence & & & \\
\rowcolor{alignmentblue}
B9.4 ตระหนักถึงความส าคัญของพฤติกรรมสุขภาพ เพศ วัย เชื้อชาติ อาชีพ เศรษฐกิจ สังคม วัฒนธรรมและ สิ่งแวดล้อม ตลอดจนปัจจัยอื่นๆ ที่มีผลกระทบต่อการบ าบัดรักษา ฟื้นฟูสภาพ สร้างเสริมสุขภาพและป้องกัน & & & \\
\rowcolor{alignmentblue}
โรค กลุ่มอาการและภาวะผิดปกติในข้อ B9.2 ซึ่งรวมถึง & & & \\
\rowcolor{alignmentblue}
B9.4.1 Emotional and behavioral factors (e.g., drug-induced interstitial nephritis, diet) & & & \\
\rowcolor{alignmentblue}
B9.4.2 Influence on person, family, and society (e.g., hemodialysis, living related kidney donation, transplants) & & & \\
\rowcolor{alignmentblue}
B9.4.3 Occupational and other environmental risk factors (e.g., heavy metals) & & & \\
\rowcolor{alignmentblue}
B9.4.4 Gender and ethnic factors (e.g., disease progression, urinary tract infections) & & & \\
\rowcolor{alignmentblue}
\textbf{B10 Reproductive System} & \textbf{10} & \textbf{30} & TBD \\
\rowcolor{alignmentblue}
B10.1 มีความรู้เรื่องภาวะปกติของระบบสืบพันธุ์และภาวะปกติในระยะปริก าเนิดดังต่อไปนี้ & & & \\
\rowcolor{alignmentblue}
B10.1.1 Embryonic development, fetal growth and development, and perinatal changes & & & \\
\rowcolor{alignmentblue}
B10.1.2 Organ structure and function & & & \\
\rowcolor{alignmentblue}
B10.1.2.1 female structure, including breast & & & \\
\rowcolor{alignmentblue}
B10.1.2.2 female function - menstrual cycle, puberty, menopause & & & \\
\rowcolor{alignmentblue}
B10.1.2.3 male structure & & & \\
\rowcolor{alignmentblue}
B10.1.2.4 male function - spermatogenesis, puberty & & & \\
\rowcolor{alignmentblue}
B10.1.2.5 intercourse, orgasm & & & \\
\rowcolor{alignmentblue}
B10.1.2.6 pregnancy, including labor and delivery, the puerperium, lactation, gestational uterus, placenta & & & \\
\rowcolor{alignmentblue}
B10.1.3 Cell/tissue structure and function including hypothalamic-pituitary-gonadal axis, sex steroids, and gestational hormones & & & \\
\rowcolor{alignmentblue}
B10.1.4 Reproductive system defense mechanisms and normal flora & & & \\
\rowcolor{alignmentblue}
B10.1.5 Repair, regeneration, and changes associated with stage of life & & & \\
\rowcolor{alignmentblue}
B10.2 มีความรู้เรื่องสาเหตุ พยาธิก าเนิด พยาธิสรีรวิทยา พยาธิสภาพ ตลอดจนหลักการ วินิจฉัยสาเหตุของการเกิดโรค กลุ่มอาการและภาวะผิดปกติทางระบบสืบพันธุ์และระยะปริก าเนิดดังต่อไปนี้ & & & \\
\rowcolor{alignmentblue}
B10.2.1 Hereditary, congenital and structural disorders & & & \\
\rowcolor{alignmentyellow}
B10.2.1.1 กลุ่มโรคและกลุ่มอาการที่ต้องรู้ & & & \\
\rowcolor{alignmentred}
B10.2.1.2 กลุ่มโรคและกลุ่มอาการที่ควรรู้ & & & \\
\rowcolor{alignmentred}
(1) imperforate hymen & & & \\
\rowcolor{alignmentred}
(2) hypospadias & & & \\
\rowcolor{alignmentred}
(3) undescended testis & & & \\
\rowcolor{alignmentred}
(4) ambiguous genitalia & & & \\
\rowcolor{alignmentblue}
B10.2.2 Infectious, inflammatory, and immunologic disorders & & & \\
\rowcolor{alignmentyellow}
B10.2.2.1 กลุ่มโรคและกลุ่มอาการที่ต้องรู้ & & & \\
\rowcolor{alignmentred}
B10.2.2.2 กลุ่มโรคและกลุ่มอาการที่ควรรู้ & & & \\
\rowcolor{alignmentyellow}
(1) sexually transmitted disease (syphilis, gonococcal infection, chancroid, non-specific urethritis, condyloma acuminata, lymphogranuloma venerum, trichomoniasis) & & & \\
\rowcolor{alignmentred}
(1) Intrauterine and perinatal infections (rubella, syphilis, gonococcal infection, tetanus neonatorum, sepsis, HIV, chlamydial infection) & & & \\
\rowcolor{alignmentyellow}
(2) Bartholinitis \& abscess, Bartholin cyst & & & \\
\rowcolor{alignmentyellow}
(3) vulvovaginitis & & & \\
\rowcolor{alignmentyellow}
(4) cervicitis & & & \\
\rowcolor{alignmentyellow}
(5) pelvic inflammatory disease (endometritis, salpingitis, oophoritis, tubo-ovarian abscess) & & & \\
\rowcolor{alignmentyellow}
(6) orchitis and epididymitis & & & \\
\rowcolor{alignmentyellow}
(7) urethritis & & & \\
\rowcolor{alignmentyellow}
(8) phimosis and paraphimosis & & & \\
\rowcolor{alignmentyellow}
(9) mastitis, breast abscess & & & \\
\rowcolor{alignmentyellow}
(10) perinatal infections (hepatitis B virus) & & & \\
\rowcolor{alignmentblue}
B10.2.3 Traumatic and mechanical disorders & & & \\
\rowcolor{alignmentyellow}
B10.2.3.1 กลุ่มโรคและกลุ่มอาการที่ต้องรู้ & & & \\
\rowcolor{alignmentred}
B10.2.3.2 กลุ่มโรคและกลุ่มอาการที่ควรรู้ & & & \\
\rowcolor{alignmentyellow}
(1) stress incontinence & & & \\
\rowcolor{alignmentred}
(1) hydrocoele of testis, spermatocoele & & & \\
\rowcolor{alignmentyellow}
(2) foreign body & & & \\
\rowcolor{alignmentred}
(2) torsion of testis & & & \\
\rowcolor{alignmentred}
(3) female genital prolapse (cystocoele, rectocoele, prolapse uterus) & & & \\
\rowcolor{alignmentred}
(4) fistula involving female genital tracts & & & \\
\rowcolor{alignmentblue}
B10.2.4 Neoplastic disorders and tumor-like condition & & & \\
\rowcolor{alignmentyellow}
B10.2.4.1 กลุ่มโรคและกลุ่มอาการที่ต้องรู้ & & & \\
\rowcolor{alignmentred}
B10.2.4.2 กลุ่มโรคและกลุ่มอาการที่ควรรู้ & & & \\
\rowcolor{alignmentyellow}
(1) hyperplasia of prostate & & & \\
\rowcolor{alignmentred}
(1) hydatidiform mole & & & \\
\rowcolor{alignmentyellow}
(2) cervical polyp & & & \\
\rowcolor{alignmentred}
(2) choriocarcinoma & & & \\
\rowcolor{alignmentred}
(3) benign and malignant neoplasm of breast, vulva, uterus, cervix, ovary, placenta, prostate gland and testes & & & \\
\rowcolor{alignmentblue}
B10.2.5 Metabolic and regulatory disorders & & & \\
\rowcolor{alignmentyellow}
B10.2.5.1 กลุ่มโรคและกลุ่มอาการที่ต้องรู้ & & & \\
\rowcolor{alignmentred}
B10.2.5.2 กลุ่มโรคและกลุ่มอาการที่ควรรู้ & & & \\
\rowcolor{alignmentyellow}
(1) menstrual disorders & & & \\
\rowcolor{alignmentred}
(1) endometriosis & & & \\
\rowcolor{alignmentyellow}
(2) menopausal syndrome & & & \\
\rowcolor{alignmentyellow}
(3) abnormal uterine or vaginal bleeding & & & \\
\rowcolor{alignmentblue}
B10.2.6 Disorders relating to pregnancy, the puerperium, and the post partum period & & & \\
\rowcolor{alignmentyellow}
B10.2.7.1 กลุ่มโรคและกลุ่มอาการที่ต้องรู้ & & & \\
\rowcolor{alignmentred}
B10.2.7.2 กลุ่มโรคและกลุ่มอาการที่ควรรู้ & & & \\
\rowcolor{alignmentyellow}
(1) abortion & & & \\
\rowcolor{alignmentred}
(1) ectopic pregnancy & & & \\
\rowcolor{alignmentyellow}
(2) hyperemesis gravidarum & & & \\
\rowcolor{alignmentred}
(2) multiple gestation & & & \\
\rowcolor{alignmentyellow}
(3) breast infection associated with childbirth & & & \\
\rowcolor{alignmentred}
(3) hypertensive disorder in pregnancy & & & \\
\rowcolor{alignmentyellow}
(4) disorders of breast and lactation associated with childbirth & & & \\
\rowcolor{alignmentred}
(4) septic abortion & & & \\
\rowcolor{alignmentred}
(5) dead fetus in utero & & & \\
\rowcolor{alignmentyellow}
(5) uterine rupture & & & \\
\rowcolor{alignmentred}
(6) antepartum hemorrhage (placenta previa, abruption placentae) & & & \\
\rowcolor{alignmentred}
(7) chorio-amnionitis & & & \\
\rowcolor{alignmentred}
(8) postpartum hemorrhage & & & \\
\rowcolor{alignmentred}
(10) amniotic fluid aspiration & & & \\
\rowcolor{alignmentred}
(11) puerperal infection & & & \\
\rowcolor{alignmentblue}
B10.2.7 Disorders in the perinatal period & & & \\
\rowcolor{alignmentyellow}
B10.2.8.1 กลุ่มโรคและกลุ่มอาการที่ต้องรู้ & & & \\
\rowcolor{alignmentred}
B10.2.8.2 กลุ่มโรคและกลุ่มอาการที่ควรรู้ & & & \\
\rowcolor{alignmentyellow}
(1) perinatal jaundice & & & \\
\rowcolor{alignmentred}
(1) respiratory distress in newborn & & & \\
\rowcolor{alignmentblue}
B10.3 มีความรู้เรื่องหลักการรักษาและหลักการใช้ยา โดยเฉพาะยาในบัญชียาหลักแห่งชาติของโรค กลุ่มอาการ & & & \\
\rowcolor{alignmentblue}
และภาวะผิดปกติในข้อ B10.2 และรวมถึงยาในกลุ่มต่อไปนี้ & & & \\
\rowcolor{alignmentblue}
B10.3.1 Prostaglandins, prostaglandin antagonists and oxytocics & & & \\
\rowcolor{alignmentblue}
B10.3.2 Myometrial relaxants & & & \\
\rowcolor{alignmentblue}
B10.3.4 Contraceptives & & & \\
\rowcolor{alignmentblue}
B10.3.5 Other methods of contraception (e.g., condoms) & & & \\
\rowcolor{alignmentblue}
B10.3.6 Estrogen, progestogen replacement and treatment of menopause & & & \\
\rowcolor{alignmentblue}
B10.3.9 Stimulators and inhibitors of lactation & & & \\
\rowcolor{alignmentblue}
B10.3.10 Male sex hormones & & & \\
\rowcolor{alignmentblue}
B10.4 ตระหนักถึงความส าคัญของพฤติกรรมสุขภาพ เพศ วัย เชื้อชาติ อาชีพ เศรษฐกิจ สังคม วัฒนธรรมและ สิ่งแวดล้อม ตลอดจนปัจจัยอื่นๆ ที่มีผลกระทบต่อการบ าบัดรักษา ฟื้นฟูสภาพ สร้างเสริมสุขภาพและ & & & \\
\rowcolor{alignmentblue}
ป้องกันโรค กลุ่มอาการและภาวะผิดปกติในข้อ B10.2 ซึ่งรวมถึง & & & \\
\rowcolor{alignmentblue}
B10.4.1 Emotional and behavioral factors (e.g., sexually transmitted diseases) & & & \\
\rowcolor{alignmentblue}
B10.4.2 Influence on person, family, and society (e.g., infertility) & & & \\
\rowcolor{alignmentblue}
B10.4.3 Occupational and other environmental risk factors (e.g., radiation) & & & \\
\rowcolor{alignmentblue}
B10.4.4 Family planning and pregnancy (e.g., unwanted, teenage pregnancy) & & & \\
\rowcolor{alignmentblue}
B10.4.5 Gender identity, sexual orientation, sexuality, libido & & & \\
\rowcolor{alignmentblue}
B10.4.6 Effects of traumatic stress syndrome, violence, rape, child abuse & & & \\
\rowcolor{alignmentblue}
\textbf{B11 Endocrine System} & \textbf{10} & \textbf{30} & TBD \\
\rowcolor{alignmentblue}
B11.1 มีความรู้เรื่องภาวะปกติของระบบต่อมไร้ท่อดังต่อไปนี้ & & & \\
\rowcolor{alignmentblue}
B11.1.1 Embryonic development, fetal maturation, and perinatal changes & & & \\
\rowcolor{alignmentblue}
B11.1.2 Organ structure and function & & & \\
\rowcolor{alignmentblue}
B11.1.2.1. hypothalamus, posterior and anterior pituitary gland & & & \\
\rowcolor{alignmentblue}
B11.1.2.2. thyroid gland & & & \\
\rowcolor{alignmentblue}
B11.1.2.3. parathyroid glands & & & \\
\rowcolor{alignmentblue}
B11.1.2.4. adrenal cortex, adrenal medulla & & & \\
\rowcolor{alignmentblue}
B11.1.2.5. pancreatic islets & & & \\
\rowcolor{alignmentblue}
B11.1.2.6. ovary and testis & & & \\
\rowcolor{alignmentblue}
B11.1.2.7. adipose tissue & & & \\
\rowcolor{alignmentblue}
B11.1.3 Cell/tissue structure and function, including hormone synthesis, secretion, action, and metabolism & & & \\
\rowcolor{alignmentblue}
B11.1.3.1 peptide hormones & & & \\
\rowcolor{alignmentblue}
B11.1.3.2 steroid hormones, including vitamin D & & & \\
\rowcolor{alignmentblue}
B11.1.3.3 thyroid hormones & & & \\
\rowcolor{alignmentblue}
B11.1.3.4 catecholamine hormones & & & \\
\rowcolor{alignmentblue}
B11.1.3.5 renin-angiotensin system & & & \\
\rowcolor{alignmentblue}
B11.1.4 Repair, regeneration, and changes associated with stage of life & & & \\
\rowcolor{alignmentblue}
B11.2 มีความรู้เรื่อง สาเหตุ พยาธิก าเนิด พยาธิสรีรวิทยา พยาธิสภาพ ตลอดจนหลักการวินิจฉัยสาเหตุของการ เกิดโรค กลุ่มอาการ หรือภาวะผิดปกติทางระบบต่อมไร้ท่อ ดังต่อไปนี้ & & & \\
\rowcolor{alignmentblue}
B11.2.1 Hereditary, congenital and structural disorders & & & \\
\rowcolor{alignmentyellow}
B11.2.1.1 กลุ่มโรคและกลุ่มอาการที่ต้องรู้ & & & \\
\rowcolor{alignmentred}
B11.2.1.2 กลุ่มโรคและกลุ่มอาการที่ควรรู้ & & & \\
\rowcolor{alignmentred}
(1) ectopic thyroid & & & \\
\rowcolor{alignmentblue}
B11.2.2 Infectious, inflammatory, and immunologic disorders & & & \\
\rowcolor{alignmentyellow}
B11.2.2.1 กลุ่มโรคและกลุ่มอาการที่ต้องรู้ & & & \\
\rowcolor{alignmentred}
B11.2.2.2 กลุ่มโรคและกลุ่มอาการที่ควรรู้ & & & \\
\rowcolor{alignmentyellow}
(1) thyrotoxicosis (Graves’ disease) & & & \\
\rowcolor{alignmentred}
(1) thyroiditis & & & \\
\rowcolor{alignmentyellow}
(2) toxic adenoma & & & \\
\rowcolor{alignmentblue}
B11.2.3. Neoplastic disorders and tumor-like conditions & & & \\
\rowcolor{alignmentyellow}
B11.2.3.1 กลุ่มโรคและกลุ่มอาการที่ต้องรู้ & & & \\
\rowcolor{alignmentred}
B11.2.3.2 กลุ่มโรคและกลุ่มอาการที่ควรรู้ & & & \\
\rowcolor{alignmentred}
(1) benign and malignant neoplasm of endocrine gland & & & \\
\rowcolor{alignmentblue}
B11.2.4 Metabolic and regulatory disorders & & & \\
\rowcolor{alignmentyellow}
B11.2.4.1 กลุ่มโรคและกลุ่มอาการที่ต้องรู้ & & & \\
\rowcolor{alignmentred}
B11.2.4.2 กลุ่มโรคและกลุ่มอาการที่ควรรู้ & & & \\
\rowcolor{alignmentyellow}
(1) diabetes mellitus, hyperglycemic crisis & & & \\
\rowcolor{alignmentred}
(1) pituitary, hypothalamus, parathyroid, pancreatic islet disorders, adrenal disorders) & & & \\
\rowcolor{alignmentyellow}
(2) hypoglycemia & & & \\
\rowcolor{alignmentyellow}
(3) thyroid disorders; goiter, iodine deficiency, hypothyroidism & & & \\
\rowcolor{alignmentred}
(2) Cushing’s syndrome & & & \\
\rowcolor{alignmentred}
(3) Addison’s disease, pituitary gland disorders, diabetes insipidus, parathyroid gland disorders & & & \\
\rowcolor{alignmentyellow}
(4) metabolic syndrome & & & \\
\rowcolor{alignmentyellow}
(5) thyroid crisis or storm & & & \\
\rowcolor{alignmentred}
(4) cerebral salt wasting & & & \\
\rowcolor{alignmentred}
(5) syndrome of inappropriate secretion of antidiuretic hormone (SIADH) & & & \\
\rowcolor{alignmentblue}
B11.2.5 idiopathic disorders & & & \\
\rowcolor{alignmentyellow}
B11.2.5.1 กลุ่มโรคและกลุ่มอาการที่ต้องรู้ & & & \\
\rowcolor{alignmentred}
B11.2.5.2 กลุ่มโรคและกลุ่มอาการที่ควรรู้ & & & \\
\rowcolor{alignmentred}
(1) hirsutism & & & \\
\rowcolor{alignmentblue}
B11.2.6 Others & & & \\
\rowcolor{alignmentyellow}
B11.2.6.1 กลุ่มโรคและกลุ่มอาการที่ต้องรู้ & & & \\
\rowcolor{alignmentred}
B11.2.6.2 กลุ่มโรคและกลุ่มอาการที่ควรรู้ & & & \\
\rowcolor{alignmentyellow}
(1) disorders of lipoprotein metabolism and lipidemia & & & \\
\rowcolor{alignmentblue}
B11.3 มีความรู้เรื่องหลักการรักษาและหลักการใช้ยา โดยเฉพาะยาในบัญชียาหลักแห่งชาติของโรค กลุ่มอาการและ & & & \\
\rowcolor{alignmentblue}
ภาวะผิดปกติในข้อ B11.2 และรวมถึงยาในกลุ่มต่อไปนี้ & & & \\
\rowcolor{alignmentblue}
B11.3.1 Drugs used in diabetes (insulins and oral antidiabetic drugs) & & & \\
\rowcolor{alignmentblue}
B11.3.2 Thyroid hormones and antithyroid drugs & & & \\
\rowcolor{alignmentblue}
B11.3.3 Corticosteroids & & & \\
\rowcolor{alignmentblue}
B11.3.4 Hypothalamic and pituitary hormones (including posterior pituitary hormones antagonists) & & & \\
\rowcolor{alignmentblue}
B11.3.5 Drugs affecting bone metabolism & & & \\
\rowcolor{alignmentblue}
B11.3.6 Other endocrine drugs (bromocriptine, other dopaminergic drugs and drugs affecting gonadotrophins) & & & \\
\rowcolor{alignmentblue}
B11.3.7 Lipid lowering agents & & & \\
\rowcolor{alignmentblue}
B11.4 ตระหนักถึงความส าคัญของพฤติกรรมสุขภาพ เพศ วัย เชื้อชาติ อาชีพ เศรษฐกิจ สังคม วัฒนธรรมแล สิ่งแวดล้อม ตลอดจนปัจจัยอื่นๆ ที่มีผลกระทบต่อการบ าบัดรักษา ฟื้นฟูสภาพ สร้างเสริมสุขภาพและ & & & \\
\rowcolor{alignmentblue}
ป้องกันโรค/กลุ่มอาการ/ภาวะในข้อ B11.2 ซึ่งรวมถึง & & & \\
\rowcolor{alignmentblue}
B11.4.1 Emotional and behavioral factors (e.g., compliance in diabetes mellitus, factitious use of insulin psychogenic polydipsia) & & & \\
\rowcolor{alignmentblue}
B11.4.2 Influence on person, family, and society & & & \\
\rowcolor{alignmentblue}
B11.4.3 Occupational and other environmental risk factors (e.g., radiation exposure, iodine deficiency) & & & \\
\rowcolor{alignmentblue}
B11.4.4 Gender and ethnic factors & & & \\